% Modernised reading — fonds Alexandre Grothendieck, shelfmark 138, pages 1-132.
%
% This is an INTERPRETATION, not a transcription. Everything the critical
% apparatus carried moves into footnotes. In any dispute of substance the
% transcription governs: whoever cites Grothendieck must cite the
% batch-NN.fr.tex files of this folder.
%
% Pass: Opus 5.5 (claude-opus-5-5), 2026-10-10 — first pass, unchecked against
% the pages by a human. The folder was read whole: seven batches, pages 1-132,
% all transcribed. Pages 1-2 (a printed leaflet and order form of a review),
% 23, 42, 44, 47, 49, 51, 55, 66, 68, 78, 87, 91 and 132 carry nothing of his
% and are absent from the transcription; pages 3, 56, 64, 69, 75, 96, 112 and
% 120 are cover sheets whose words become headings here.
%
% What the folder is. The inventory files it in the group « Complexe de De
% Rham » (136-139), but its own title is « Cartes. Étude arithmétique » and
% nothing in it concerns a de Rham complex. It is the earliest working file we
% have of what Grothendieck later called dessins d'enfants: maps on surfaces,
% the covers of P^1 minus {0,1,infty} they define, and the action of
% Gal(Qbar/Q) on them. Stations: a Galois dictionary of descent data
% (pp. 3-13); P^1 minus points over Q, towers to study (pp. 14-22); the Fermat
% tower as the universal abelian cover, with its fibres over the seven special
% points (pp. 24-54); automorphisms of groupoids (pp. 56-62); finite subgroups
% of PGL_2 (pp. 64-67); a group-theoretic paradigm for surfaces (pp. 69-74);
% « special » profinite groups and the outer Galois action (pp. 79-95);
% pre-maps, pinned maps and their moduli schemes, with examples (pp. 96-111);
% a census of spherical maps of degree <= 6 (pp. 112-119); « operations » on
% maps (pp. 120-131).
%
% Notation fixed for the whole document. Gamma = Gal(kbar/k), from p. 14 on
% Gal(Qbar/Q), Qbar the algebraic closure of Q in C. pp. 3-13: script G,
% gothic G, script E, G, Z, gothic z as on the pages. Fermat (pp. 24-54):
% C_n, Pi_n = mu_n^3/mu_n, T = lim mu_n (= Zhat(1)), chi the cyclotomic
% character, E_n, E_infty, E-tilde WITH their index: they are the groups
% script-G' of p. 6 for the pair C_n -> C_1, not the script E of p. 5. K(g)
% the T-torsor of compatible roots of g. Profinite (pp. 79-90): Sigma_nu,
% L_nu, Pi_nu = Lhat_nu; A-tilde_nu, A_nu, A°_nu; Gamma-tilde_nu, Gamma_nu,
% Gamma°_nu their quotients by Pi_nu. What pp. 83-86 write as a curly
% T-tilde_nu (and T_2, T^!_nu) is written Gamma-tilde_nu (Gamma_2,
% Gamma^!_nu) here; the curly T_nu is kept for the DISCRETE Teichmüller group
% B_nu/L_nu of p. 89 only. Maps: 0 -> vertices, 1 -> edge midpoints,
% infty -> face centres; C a map, C_0 a pre-map, gamma_{C_0} = Aut(C_0);
% script C_2^+ the oriented cartographic group; script C (pp. 92-95) the
% torsor of algebraic closures — three different C's, kept apart.
%
% Substantive departures from the pages, each footnoted where it occurs:
% — p. 5-7: descent data are CONTINUOUS sections; « de prés. finie sur k »
%   read as over kbar; p. 10 « invariant » vs « centralise » is an
%   implication, not an equivalence (equivalence after shrinking Gamma_i,
%   or when the centre is trivial).
% — p. 14: G = Gamma.Qbar^* is the group of automorphisms fixing 0 and infty;
%   the full group also has t -> 1/t.
% — p. 18: Q[R'] = Q(mu_n, 2^{1/n}), not Q(2^{1/n}).
% — p. 20: Kummer theory has coefficients mu_n, not Z/n; lim Q^*/Q^*n only
%   injects into {±1} x Zhat^P.
% — p. 27 (13): Aut_Q(C_n) = E_n only for n >= 4, or as automorphisms of the
%   pair C_n -> C_1.
% — p. 29 (20) and p. 33 (30): Pi_infty acts on p^{-1}(i) and p_3^{-1}(omega)
%   through multiplication by 2, resp. 3; (20) reads with 2Z_i.
% — p. 30 (23): with t = -X/Z of (3''), the formulas describe the fibre over
%   -jbar; the page's name alpha = -j is kept, the swap is harmless.
% — p. 35 (40): Gamma -> Aff(1, Zhat) has index 2 (sqrt 2 in Q(mu_8)).
% — p. 40: the second line of alpha'_n(infty) conflicts with (43); first line
%   kept.
% — p. 50: the parametrisation is of x^2+y^2 = z^2; over Q(i) equivalent.
% — p. 57: the product in b) runs over x != a_i; pp. 59-60: Aut C -> S_{C_0}
%   surjective for C connected; p. 60 corollaries need C' a groupoid.
% — p. 67: the Sylow subgroup is (Z/p)^r, and p | N is outside the tame
%   hypothesis of the heading.
% — p. 80: Pi -> A°_nu injective needs nu >= 2.
% — p. 83: Gamma -> S_{nu+1} factors through (Z/(nu+1))^*, not Z/(nu+1).
% — p. 84: t_nu as written does not preserve the relation; a corrected
%   formula is given.
% — p. 89: T°_nu is the pure mapping class group = pure spherical braids
%   modulo their centre Z/2.
% — p. 97-99: co-pinned maps are rigid in genus 0 but NOT in general (torus,
%   one vertex, two edges, one face); the count of p. 98 is a mass formula.
% — p. 102: the action of Gamma on C^alpha/gamma is non-trivial for SOME C_0.
% — p. 110: roots of 25c^2-32c+16 are 16/25 ± (12/25)i, not ±(24/25)i.
% — p. 122: the H_alpha must be normal for G_alpha to be a group.
% — p. 129-130: the group of case II is R_{>0}, not C^*; the oranges of every
%   order n have z -> az (a > 0); the conclusion of p. 130 survives.
%
% Doubts flagged, not resolved: the claim of p. 101 that for H = 1 the oriented
% type is not determined by the cover; the group Aff(1,C) given for case III
% on p. 129; the overall reading of pp. 52-53, 76 and 125 (heavily corrected
% pages).
%
% Modern names given and footnoted as ours: Galois descent; fundamental exact
% sequence (SGA 1); Fermat curves as the maximal abelian covers; Kummer theory
% and its open image; Belyi function; Klein's classification; Riemann-Hurwitz;
% Dehn-Nielsen-Baer; ends of surfaces; Earle-Eells; mapping class group and
% M_{0,n}; congruence subgroup property (Asada); outer Galois representation,
% Grothendieck-Teichmüller group (Drinfeld, Ihara); Belyi's theorem; triangle
% (orbifold) groups; Riemann existence theorem; Hurwitz schemes (Fulton);
% oriented maps as permutation pairs (Jones-Singerman 1978); Shabat polynomial.
%
% What the folder announces and never establishes: the transport bijections of
% p. 39-40 (stopped after one of the two needed); the « problems » of p. 17;
% the isomorphisms Gamma^!_2 = Gamma_2^{S_3} = Gamma^!_nu wished for on p. 86;
% the questions of p. 89; the polynomial P(c) of pp. 105-111 and whether its
% roots avoid delta(c); the « exceptional case » of p. 127; the conformal
% argument of p. 130; the proof of the theorem of p. 125 beyond its first step;
% the conditions a)-d) of p. 131, whose comparison is not made.
\documentclass[11pt,a4paper]{article}
\input{../preamble/grothendieck}

\folder{138}
\pages{1}{132}
\dating{[vers 1976-1978]}
\watermark{Édition de démonstration\\interprétation personnelle de l'œuvre}
\foldertitle{Cartes. Etude arithmétique (1976, 77 ?) : bon de commande (s.d.), notes manuscrites (s.d.) — lecture modernisée du dossier entier}

\begin{document}

\section*{Résumé}

\begin{resume}
Dessinez sur une sphère un petit graphe : quelques points, quelques traits qui
les joignent sans se croiser. La sphère se trouve découpée en régions. Ce
dessin anodin --- une \emph{carte} --- se révèle porter une structure rigide :
il existe une manière, et essentiellement une seule, de voir la sphère comme
la droite des nombres complexes complétée par un point à l'infini, de sorte que
le dessin soit l'image réciproque du segment $[0,1]$ par une fonction
rationnelle qui ne « se replie » qu'au-dessus de $0$, de $1$ et de l'infini. Les
coefficients de cette fonction sont alors des nombres algébriques. Le groupe
des symétries de tous les nombres algébriques --- le groupe de Galois absolu
de $\mathbb{Q}$, l'objet le plus mystérieux de l'arithmétique --- agit donc sur
les dessins : il change un dessin en un autre. Ces feuillets, écrits vers
1976--1978, sont le premier chantier que nous ayons de ce que Grothendieck
appellera plus tard les « dessins d'enfants ».

Le problème est de comprendre cette action. Elle est fidèle à grande échelle
mais invisible à petite : sur un dessin très simple, elle ne fait presque
rien, et toute la difficulté est d'en trouver des traces calculables. Le
dossier attaque de plusieurs côtés à la fois. Il commence par un
\emph{dictionnaire} : « descendre » un objet géométrique d'un corps vers un
corps plus petit revient à scinder une suite exacte de groupes, et les
obstructions vivent dans la cohomologie. Il étudie ensuite à fond le cas où
tout est abélien : la tour des courbes de Fermat $x^n + y^n + z^n = 0$, qui
recouvre la sphère privée de trois points comme un pavage par $2n^2$
triangles, et où l'action de Galois s'écrit entièrement en termes de racines
de l'unité et de racines $n$-ièmes de $2$. Il change ensuite d'altitude : une
surface trouée est connue, à déformation près, par son groupe fondamental muni
de la donnée des petits lacets autour des trous ; le groupe de Galois agit sur
la version « profinie » de ce groupe, et Grothendieck nomme les groupes dans
lesquels cette action prend ses valeurs. Viennent enfin les objets eux-mêmes :
des cartes « épinglées », dont les espaces de modules sont des ensembles finis
de points algébriques ; un recensement à la main de toutes les cartes de la
sphère jusqu'à six arêtes ; le calcul explicite de la fonction rationnelle
attachée à quelques arbres ; et l'idée d'« opérations » sur les cartes (passer
au dual, subdiviser) qui commutent avec tout ce qui précède.

Une analogie aide à tenir le fil. Un polynôme à coefficients rationnels a des
racines que Galois permute ; on ne voit pas ces permutations en regardant les
racines une à une, mais on les voit dans les relations qu'elles satisfont.
Ici les « racines » sont des dessins, et les « relations » sont les manières
dont un dessin se déduit d'un autre --- en le recouvrant, en le
symétrisant, en lui appliquant une opération. Plus on accumule de ces
relations, plus l'action de Galois est contrainte.

Le dossier est un atelier, non un exposé : les pages s'interrompent, se
reprennent, posent des questions qu'elles ne tranchent pas, et quelques
énoncés y sont faux tels qu'écrits --- la présente lecture les corrige et le
dit en note. Mais presque tout ce qui fera la célébrité de l'\emph{Esquisse
d'un programme} (1984) y est déjà en germe : le groupe cartographique,
l'action extérieure de Galois sur le groupe fondamental de la droite privée de
trois points, la tour des groupes de Teichmüller, le rôle exceptionnel du
groupe des permutations de $\{0,1,\infty\}$. Les noms modernes à chercher
ensuite sont ceux de Belyi, de Drinfeld et d'Ihara.

\keywords{dessins d'enfants, Belyi map, combinatorial map, cartographic group, Galois descent, fundamental exact sequence, Fermat curve, Kummer theory, cyclotomic character, outer Galois representation, profinite fundamental group, Grothendieck-Teichmüller group, mapping class group, Dehn-Nielsen-Baer theorem, pure braid group, moduli space of pointed rational curves, finite subgroups of PGL2, Riemann-Hurwitz formula, groupoid automorphisms, Hurwitz space, Riemann existence theorem, Shabat polynomial}
\end{resume}

\section*{Le fil du dossier, et les conventions}

\begin{enumerate}
\item[1.] \emph{Dictionnaire galoisien des descentes} (pages 3 à 13) : données de
descente comme scindages, et l'extension centrale qui mesure la descente d'un
revêtement avec ses automorphismes.
\item[2.] \emph{Sur $\mathbb{Q}$, la droite projective privée de points}
(pages 14 à 22) : les cas $\operatorname{Card} S = 0, 1, 2, 3, \geqslant 4$, la
liste des tours à étudier, et des calculs sur $z \mapsto z^{n}$, la fonction
$S_3$-invariante $g$ et la théorie de Kummer sur $\mathbb{Q}$.
\item[3.] \emph{La tour de Fermat} (pages 24 à 40) : le revêtement abélien
universel de la droite privée de $0, 1, \infty$, ses fibres au-dessus des sept
points spéciaux, et leur description comme ensembles à opérateurs.
\item[4.] \emph{Fermat, suite} (pages 41 à 54) : la version fonctorielle en un
ensemble à trois éléments, l'octaèdre, la tour diédrale, les endomorphismes de
la conique.
\item[5.] \emph{Automorphismes de groupoïdes} (pages 56 à 62), et le graphe
$K(I)$ à deux sommets.
\item[6.] \emph{La double pyramide} (page 63) : les triangles d'une carte et
les rotations autour des sommets.
\item[7.] \emph{Sous-groupes finis de $PGL_2$, cas modéré} (pages 64 à 67).
\item[8.] \emph{Paradigme « gruppentheoretisch » des surfaces} (pages 69 à 74).
\item[9.] \emph{Quelques problèmes spéciaux} (pages 75 à 77) : les drapeaux d'une
carte.
\item[10.] \emph{Groupes profinis spéciaux et action de Galois} (pages 79 à 90).
\item[11.] \emph{Le torseur des clôtures algébriques} (pages 92 à 95).
\item[12.] \emph{Précartes, cartes co-épinglées, leurs schémas de modules, et des
exemples} (pages 96 à 111).
\item[13.] \emph{Cartes sphériques en bas degrés} (pages 112 à 119).
\item[14.] \emph{« Opérations » sphériques sur les cartes} (pages 120 à 131).
\end{enumerate}

Les stations 1 et 10--11 sont la même idée à deux altitudes : la page 6
décrit les automorphismes d'un revêtement par une suite exacte au-dessus de
$\Gamma$, la page 82 décrit l'action de $\Gamma$ sur le groupe fondamental tout
entier. La tour de Fermat (stations 2--4) est le cas « abélien » du
dictionnaire, entièrement calculé. Les stations 6, 9 et 12--14 sont la théorie
combinatoire des cartes dont tout le reste parle.

\emph{Conventions.} $k$ est un corps, $\bar{k}$ une clôture séparable,
$\Gamma = \mathrm{Gal}(\bar{k}/k)$ ; à partir de la page 14, $k = \mathbb{Q}$,
$\bar{\mathbb{Q}}$ est la clôture algébrique de $\mathbb{Q}$ dans $\mathbb{C}$ et
$\Gamma = \mathrm{Gal}(\bar{\mathbb{Q}}/\mathbb{Q})$. On note
\[
\mathrm{T} = \varprojlim_{n} \mu_{n}(\bar{\mathbb{Q}}) \;(\simeq \widehat{\mathbb{Z}}(1)),
\qquad
\chi : \Gamma \to \widehat{\mathbb{Z}}^{*} = \mathrm{Aut}(\mathrm{T})
\]
le module de Tate des racines de l'unité et le caractère cyclotomique, et
$\chi_{n} = \chi \bmod n$. Pour $g \in \bar{\mathbb{Q}}^{*}$, $K(g)$ est
l'ensemble des systèmes compatibles $(\xi_{n})$ de racines $n$-ièmes de $g$,
$\xi_{nm}^{m} = \xi_{n}$ : c'est un $\mathrm{T}$-torseur, et, si $g \in
\mathbb{Q}^{*}$, $\Gamma$ y opère de façon compatible à son action sur
$\mathrm{T}$. On écrit $A \ltimes B$ un produit semi-direct où $B$ est
distingué.\footnote{Il écrit $A \cdot_{1/2} B$, ou $A \cdot B$, un point muni
d'un « $\tfrac12$ » en indice : « produit $\tfrac12$ direct » (p.~43).}

Les lettres sont celles des pages, avec trois exceptions, chacune dite en note
au point où elle sert. Les groupes $\mathcal{E}_{n}$, $\mathcal{E}_{\infty}$,
$\widetilde{\mathcal{E}}$ de la tour de Fermat sont toujours écrits avec leur
indice ou leur tilde : ce sont des groupes $\mathcal{G}'$ au sens de la page 6,
et non le groupe $\mathcal{E}$ de la page 5. Pour les groupes profinis des
pages 79 à 90, on garde les noms de la page 80 ($\widetilde{\Gamma}_{\nu}$,
$\Gamma_{\nu}$, $\Gamma^{\circ}_{\nu}$) là où les pages 83 à 86 écrivent un
$\mathcal{T}$ cursif, et l'on réserve $\mathcal{T}_{\nu}$ au groupe de
Teichmüller \emph{discret} de la page 89. Enfin trois objets portent la lettre
C : une carte $C$ et une précarte $C_{0}$ ; le groupe cartographique orienté
$\mathcal{C}_{2}^{+}$ ; et le torseur $\mathcal{C}$ des pages 92 à 95.

Une \emph{carte} est un graphe fini $K$ tracé sur une surface compacte
orientée $X$, de sorte que les composantes de $X \smallsetminus K$ soient des
disques ; ses \emph{brins} (arêtes orientées issues d'un sommet) sont permutés
par la rotation autour des sommets et par le retournement des arêtes. Dans la
correspondance avec les revêtements de la droite projective, \emph{les sommets
sont au-dessus de $0$, les milieux d'arêtes au-dessus de $1$, les centres des
faces au-dessus de $\infty$}.\footnote{C'est la convention des pages 96 et
suivantes (« $f^{-1}(0) \simeq S$, $f^{-1}(1) \simeq A$, $f^{-1}(\infty) \simeq
F$ », p.~99) et des dessins de la page 103. Les pages 24 à 54 n'ont pas encore
de convention : les trois points y jouent des rôles symétriques.}

\emph{Ce que le dossier annonce et n'établit pas} : les bijections de
« transport parallèle » de la page 39, dont une seule sur les deux nécessaires
est entamée ; les problèmes de la liste de la page 17 ; les isomorphismes
canoniques souhaités à la page 86 ; les questions de la page 89 ; le polynôme
$P(c)$ des pages 105 à 111 ; le « cas exceptionnel » de la page 127 ;
l'argument de représentation conforme de la page 130 ; la démonstration du
théorème de la page 125 au-delà de son premier pas ; la comparaison des
conditions a) à d) de la page 131.

\emph{Les feuillets.} Les pages 1 et 2 sont un dépliant imprimé et un bon de
commande d'une revue, sans rien de sa main ; ils donnent leur « bon de
commande » au titre de l'inventaire. Plusieurs feuillets sont écrits au verso
de documents dactylographiés de 1977 et 1978 (une lettre, des avis de
soutenance), d'où la datation.\footnote{Ces documents ne sont pas transcrits
et ne sont pas repris ici. La datation « [vers 1976-1978] » est celle de
l'inventaire ; le titre de l'inventaire, « (1976, 77 ?) », est sans doute
celui d'une chemise.} L'inventaire range ce dossier, avec les cotes 136, 137 et
139, sous le titre de groupe « Complexe de De Rham » ; rien dans ces 132 pages
ne porte sur un complexe de De Rham.

\section*{1. Dictionnaire galoisien des descentes (pages 3 à 13)}

\pagerange{4}{5}
\subsection*{Descente et scindages (pages 4 et 5)}

Soit $\mathfrak{X}$ un schéma connexe sur $\bar{k}$, et supposons $\bar{k}$
séparablement clos dans $H^{0}(\mathfrak{X}, \mathcal{O}_{\mathfrak{X}})$. On
considère les automorphismes de $\mathfrak{X}$ comme $k$-schéma, et ceux qui sont
$\bar{k}$-linéaires :
\[
\mathcal{G} = \mathrm{Aut}_{k}(\mathfrak{X}), \qquad
\mathfrak{G} = \mathrm{Aut}_{\bar{k}}(\mathfrak{X}).
\]
Un $k$-automorphisme de $\mathfrak{X}$ agit sur $H^{0}(\mathfrak{X},
\mathcal{O}_{\mathfrak{X}})$ et y stabilise $\bar{k}$, la clôture séparable de
$k$ qu'il contient ; d'où une suite exacte
\[
1 \to \mathfrak{G} \to \mathcal{G} \to \Gamma .
\]
\footnote{L'hypothèse sur $H^{0}$ est en marge de la page 4 ; c'est elle qui
fait exister la flèche $\mathcal{G} \to \Gamma$. Sans elle, un automorphisme
de $\mathfrak{X}$ comme $k$-schéma n'a aucune raison de respecter le
sous-corps $\bar{k}$.}
Pour $k \subset k_{1} \subset \bar{k}$, de groupe $\Gamma_{1} \subset \Gamma$, le
groupe $\mathcal{G}_{1} = \mathrm{Aut}_{k_{1}}(\mathfrak{X})$ est l'image
réciproque de $\Gamma_{1}$ dans $\mathcal{G}$.

\emph{Une donnée de descente de $\mathfrak{X}$ de $\bar{k}$ à $k$ est un
scindage continu de $\mathcal{G} \to \Gamma$}, c'est-à-dire un homomorphisme
$s : \Gamma \to \mathcal{G}$ relevant l'identité, et tel que l'action de
$\Gamma$ qu'il définit soit continue ; il permet d'écrire $\mathcal{G} \simeq
\Gamma \ltimes \mathfrak{G}$. Lorsque $\mathfrak{X}$ est quasi-projectif sur
$\bar{k}$, toute donnée de descente est effective : il revient au même de se
donner un $k$-schéma $X$ et un isomorphisme $X \otimes_{k} \bar{k} \simeq
\mathfrak{X}$. De même une descente à $k_{1}$ est un $\Gamma_{1}$-scindage,
c'est-à-dire un scindage de $1 \to \mathfrak{G} \to \mathcal{G}_{1} \to
\Gamma_{1}$, et $\mathcal{G}_{1} \simeq \Gamma_{1} \ltimes \mathfrak{G}$.%
\footnote{La page dit « scindage », sans la condition de continuité. Elle est
nécessaire : un scindage abstrait de $\mathcal{G} \to \Gamma$ ne descend pas
en général, faute que l'action de chaque élément de $\Gamma$ sur un ouvert
affine se lise à un niveau fini. C'est la descente galoisienne de Weil (1956),
reprise dans SGA~1, exposé VIII ; Grothendieck la connaissait évidemment.
L'effectivité pour $\mathfrak{X}$ quasi-projectif est sur la page.}

\pagerange{5}{7}
\subsection*{Le revêtement universel et un revêtement galoisien (pages 5 à 7)}

Soit $\widetilde{\mathfrak{X}} \to \mathfrak{X}$ un revêtement universel
profini, et
\[
\mathcal{E} = \mathrm{Aut}_{k}(\widetilde{\mathfrak{X}} \to \mathfrak{X}),
\qquad
\pi_{1} = \mathrm{Aut}(\widetilde{\mathfrak{X}}/\mathfrak{X}),
\]
le groupe des $k$-automorphismes de $\widetilde{\mathfrak{X}}$ qui relèvent un
$k$-automorphisme de $\mathfrak{X}$, et le groupe fondamental profini. Comme
tout automorphisme de $\mathfrak{X}$ se relève au revêtement universel, on a
une suite exacte
\[
1 \to \pi_{1} \to \mathcal{E} \to \mathcal{G} \to 1,
\]
d'où une action extérieure $\mathcal{G} \to \mathrm{Out}(\pi_{1})$ par
automorphismes continus.\footnote{Pour $\mathfrak{X}$ défini sur $k$, c'est la
« suite exacte fondamentale » de SGA~1, exposé IX,
$1 \to \pi_{1}(X_{\bar{k}}) \to \pi_{1}(X) \to \Gamma \to 1$, élargie ici aux
automorphismes géométriques. La note de marge sur la continuité de
$\mathrm{int}(g)|\pi_{1}$ est lue avec un doute sur « continu ».}

Soit $N \subset \pi_{1}$ un sous-groupe fermé distingué,
$\mathfrak{X}' = \widetilde{\mathfrak{X}}/N$ le revêtement galoisien
correspondant, de groupe $G = \pi_{1}/N = \mathrm{Aut}(\mathfrak{X}'/\mathfrak{X})$.
Posons
\[
\mathcal{E}' = \mathrm{Norm}_{\mathcal{E}}(N) \supset \pi_{1},
\qquad
\mathcal{G}' = \mathcal{E}'/N \xrightarrow{\;\sim\;}
\mathrm{Aut}_{k}(\mathfrak{X}' \to \mathfrak{X}),
\]
le second groupe étant celui des $k$-automorphismes du \emph{couple}
$\mathfrak{X}' \to \mathfrak{X}$. On a une suite exacte
\[
1 \to G \to \mathcal{G}' \to \mathcal{G} \to \Gamma,
\]
et l'image de $\mathcal{G}' \to \mathcal{G}$ est $\mathcal{E}'/\pi_{1}$,
c'est-à-dire le stabilisateur de $N$ pour l'action de $\mathcal{G}$ sur
l'ensemble des sous-groupes fermés distingués de $\pi_{1}$. En particulier
\emph{$\mathcal{G}' \to \mathcal{G}$ est surjectif si et seulement si $N$ est
distingué dans $\mathcal{E}$}. Si $\pi_{1}$ est topologiquement de type fini,
ses sous-groupes ouverts caractéristiques forment un système fondamental de
voisinages de $1$ (il n'y a qu'un nombre fini de sous-groupes ouverts d'indice
donné, et leur intersection est caractéristique) : il y a donc beaucoup de $N$
distingués dans $\mathcal{E}$.\footnote{La page 5 numérote (12) une formule
déjà numérotée, puis reprend à (13) ; on ne garde pas sa numérotation. Le
« caractéristiques » de la page 6 est un ajout douteux, mais c'est ce que
l'argument demande.}

Soit $\mathfrak{G}' = \mathrm{Aut}_{\bar{k}}(\mathfrak{X}' \to \mathfrak{X})$ ;
c'est le noyau de $\mathcal{G}' \to \Gamma$, et l'image réciproque de
$\mathfrak{G}$, d'où une suite exacte $1 \to G \to \mathfrak{G}' \to
\mathfrak{G}$. Le « cas favorable » où $\mathcal{G}' \to \mathcal{G}$ est
surjectif équivaut à la conjonction de : a) $\mathfrak{G}' \to \mathfrak{G}$
est surjectif ; b) $\mathcal{G}'$ et $\mathcal{G}$ ont même image dans
$\Gamma$. (Si a) et b) sont vrais et $g \in \mathcal{G}$, on choisit $g' \in
\mathcal{G}'$ de même image dans $\Gamma$, et l'on relève $g\,\bar{g'}^{-1} \in
\mathfrak{G}$ par a).) Si $\mathfrak{X}$ est de présentation finie sur
$\bar{k}$, il est défini sur une extension finie de $k$, et l'image de
$\mathcal{G}$ dans $\Gamma$ est un sous-groupe ouvert ; de même pour
$\mathcal{G}'$ si $\mathfrak{X}' \to \mathfrak{X}$ est fini.\footnote{La page
écrit « de prés.\ finie sur $k$ » ; $\mathfrak{X}$ est un $\bar{k}$-schéma, et
c'est la présentation finie sur $\bar{k}$ qui donne un corps de définition
fini.}

\pagerange{8}{8}
\subsection*{Descente d'un couple (page 8)}

Une donnée de descente de $(\mathfrak{X}' \to \mathfrak{X})$ de $\bar{k}$ à $k$
est, de même, un scindage continu de $\mathcal{G}' \to \Gamma$, qui permet
d'écrire $\mathcal{G}' \simeq \Gamma \ltimes \mathfrak{G}'$ ; une descente à
$k_{1}$ est un scindage partiel au-dessus de $\Gamma_{1}$, c'est-à-dire un
scindage de $\mathcal{G}'_{1} \to \Gamma_{1}$, où
\[
\mathcal{E}_{1} = \mathrm{Aut}_{k_{1}}(\widetilde{\mathfrak{X}} \to \mathfrak{X}),
\qquad
\mathcal{E}'_{1} = \mathcal{E}' \cap \mathcal{E}_{1},
\qquad
\mathcal{G}'_{1} = \mathrm{Aut}_{k_{1}}(\mathfrak{X}' \to \mathfrak{X})
\]
sont les images réciproques de $\Gamma_{1}$ ; on écrit alors
$\mathcal{G}'_{1} \simeq \Gamma_{1} \ltimes \mathfrak{G}'$.\footnote{En marge :
« oubli : à vérifier dans les formules ! », lu avec un doute.}

\pagerange{9}{13}
\subsection*{Descendre avec tous ses automorphismes : une extension centrale
(pages 9 à 13)}

Plaçons-nous dans le cas où $N$ est distingué dans $\mathcal{E}$ et où
$\mathcal{G} \to \Gamma$ est surjectif. Alors $\mathcal{G}'$ a une suite de
composition de quotients successifs $G$, $\mathfrak{G}$, $\Gamma$ :
\[
\overbrace{G \quad \mathfrak{G} \quad \Gamma}^{\mathcal{G}'}, \qquad
\underbrace{G \quad \mathfrak{G}}_{\mathfrak{G}'}, \qquad
\underbrace{\mathfrak{G} \quad \Gamma}_{\mathcal{G}} .
\]
L'extension $\mathfrak{G}'$ de $\mathfrak{G}$ par $G$ est de nature
géométrique ; l'extension $\mathcal{G}'$ de $\Gamma$ par $\mathfrak{G}'$
exprime l'arithmétique de la situation. Le cas qui l'intéresse est celui où
$\mathfrak{G}$ et $G$ sont finis, $\Gamma$ infini et $\mathfrak{X}$ de
présentation finie.

On cherche alors une extension finie galoisienne $k_{i}$ de $k$, de groupe
$\Gamma_{i}$ (ouvert distingué dans $\Gamma$), et une extension
\[
1 \to \mathfrak{G}' \to \mathcal{G}'_{i} \to \Gamma/\Gamma_{i} \to 1
\]
dont $\mathcal{G}'$ se déduise par l'image réciproque le long de $\Gamma \to
\Gamma/\Gamma_{i}$. Se donner cela revient à se donner un scindage partiel
continu $s : \Gamma_{i} \to \mathcal{G}'$ dont l'image soit distinguée dans
$\mathcal{G}'$. Géométriquement : un modèle $X'_{i} \to X_{i}$ de
$\mathfrak{X}' \to \mathfrak{X}$ sur $k_{i}$, d'où une extension de
$k_{i}$-schémas en groupes finis étales $1 \to G_{i} \to \mathfrak{G}'_{i} \to
\mathfrak{G}_{i} \to 1$, avec la condition que $\mathfrak{G}'_{i}$ soit
\emph{constant} --- tous les automorphismes géométriques du couple sont
définis sur $k_{i}$.

Si l'image de $s$ est distinguée, elle commute à $\mathfrak{G}'$ : les deux
sous-groupes sont distingués et d'intersection triviale. La réciproque n'est
pas vraie telle quelle, mais elle l'est après avoir rétréci $\Gamma_{i}$.%
\footnote{La page dit « ce qui revient au même ». Si $s$ centralise
$\mathfrak{G}'$, ses conjugués $g s g^{-1}$ ($g \in \mathcal{G}'$) le
centralisent aussi et diffèrent de $s$ par un homomorphisme $\Gamma_{i} \to
\mathfrak{z}$ (notation ci-dessous) ; ces conjugués ne dépendent que de
l'image de $g$ dans le groupe fini $\Gamma/\Gamma_{i}$, et l'intersection des
noyaux des homomorphismes obtenus est un sous-groupe ouvert distingué sur
lequel l'image de $s$ devient distinguée. Si $\mathfrak{z} = 1$, il n'y a rien
à rétrécir.} Pour un $k_{i}$ donné il peut y avoir plusieurs solutions non
isomorphes ; l'unicité ne s'obtient qu'en passant à une extension assez
grande. Précisément, soit
\[
Z = \mathrm{Centr}_{\mathcal{G}'}(\mathfrak{G}'),
\qquad
\mathfrak{z} = Z \cap \mathfrak{G}' = \text{centre de } \mathfrak{G}' .
\]
L'image $\Gamma_{0}$ de $Z$ dans $\Gamma$ est ouverte ($Z$ est d'indice fini,
puisque $\mathfrak{G}'$ est fini), elle correspond à une extension finie
$k_{0}$, et l'on a une \emph{extension centrale}
\[
1 \to \mathfrak{z} \to Z \to \Gamma_{0} \to 1 .
\]
Les solutions relatives à $\Gamma_{i} \subset \Gamma_{0}$ sont les scindages
partiels de cette extension au-dessus de $\Gamma_{i}$.

\begin{itemize}
\item Si $\mathfrak{z} = 1$, alors $Z \simeq \Gamma_{0}$ : il existe une
solution si et seulement si $\Gamma_{i} \subset \Gamma_{0}$, et elle est
unique. Il y a donc un corps de définition canonique $k_{0}$ du couple avec
tous ses automorphismes.
\item En général, l'obstruction à l'existence au-dessus de $\Gamma_{i}$ est la
classe de l'extension restreinte dans $H^{2}(\Gamma_{i}, \mathfrak{z})$
($\Gamma_{i}$ opérant trivialement), et les solutions, quand il y en a,
forment un torseur sous $H^{1}(\Gamma_{i}, \mathfrak{z}) =
\mathrm{Hom}(\Gamma_{i}, \mathfrak{z})$. L'obstruction meurt sur un
$\Gamma_{i}$ assez petit : $Z$ étant profini et $\mathfrak{z}$ fini, un
sous-groupe ouvert de $Z$ rencontre $\mathfrak{z}$ trivialement, et son image
convient.\footnote{La page s'arrête sur « $H^{1}(\Gamma_{i}, \mathfrak{z}) =
\mathrm{Hom}(\Gamma_{i}, \mathfrak{z})$ \ldots » ; la dernière phrase est de
nous. Les groupes de cohomologie sont la cohomologie continue des groupes
profinis ; la topologie de $\mathcal{G}'$, que la page ne précise pas, est
celle pour laquelle les scindages continus sont les données de descente.}
\end{itemize}

\pagerange{14}{17}
\section*{2. Sur $\mathbb{Q}$ : la droite projective privée de points (pages 14 à 17)}

Soit $\hat{\mathfrak{X}} = \mathbb{P}^{1}_{\bar{\mathbb{Q}}}$, $S$ une partie
finie de ses points, $\mathfrak{X} = \hat{\mathfrak{X}} \smallsetminus S$, et
$\mathfrak{X}'$ le revêtement abélien maximal de $\mathfrak{X}$.\footnote{La
page écrit « $\bar{K}$ » pour $\bar{k}$ et précise « modéré » entre crochets ;
en caractéristique nulle tout revêtement est modéré.}

\textbf{A)} Si $\operatorname{Card} S \in \{0, 1\}$, $\mathfrak{X}$ est
simplement connexe et $\mathfrak{X}' = \mathfrak{X}$.

\textbf{B)} Si $\operatorname{Card} S = 2$, $\mathfrak{X}$ est canoniquement un
torseur sous une forme de $\mathbb{G}_{m}$ ; choisir l'un des points de $S$
comme $\infty$ et une origine $1$ le trivialise, et le couple
$(\hat{\mathfrak{X}}, \mathfrak{X})$ descend en $(\mathbb{P}^{1}_{\mathbb{Q}},
\mathbb{G}_{m,\mathbb{Q}})$. Le groupe des automorphismes fixant $0$ et $\infty$
est alors $\Gamma \ltimes \bar{\mathbb{Q}}^{*}$.\footnote{La page écrit
$\mathcal{G} = \Gamma \cdot \bar{\mathbb{Q}}^{*}$. Le groupe de tous les
$\mathbb{Q}$-automorphismes de $\mathbb{G}_{m,\bar{\mathbb{Q}}}$ contient de plus
$t \mapsto 1/t$, qui échange $0$ et $\infty$ : on lit $\mathcal{G}$ comme le
groupe qui respecte le choix de $\infty$ fait juste avant.} Le revêtement
d'indice $n$ est $X_{n} = \mathbb{G}_{m} \xrightarrow{t \mapsto t^{n}}
\mathbb{G}_{m}$, principal de groupe $\mu_{n}$ ; l'homomorphisme
$\mathcal{G}_{n} \to \mathcal{G}$ est l'identité sur $\Gamma$ et l'élévation à
la puissance $n$ sur $\bar{\mathbb{Q}}^{*}$. À la limite, le groupe des
automorphismes du pro-revêtement abélien maximal est $\Gamma \ltimes
\varprojlim_{n} \bar{\mathbb{Q}}^{*}$, la limite étant prise pour les
élévations aux puissances, extension de $\bar{\mathbb{Q}}^{*}$ par
$\mathrm{T}$.\footnote{La page écrit $\Gamma \cdot
T_{\infty}(\bar{\mathbb{Q}}^{*})$ ; ce n'est pas le module de Tate (qui est le
noyau $\mathrm{T}$) mais la limite tout entière.}

\textbf{C)} Si $\operatorname{Card} S = 3$, $\mathfrak{G} \simeq \mathfrak{S}_{S}
\simeq \mathfrak{S}_{3}$, et il y a une façon canonique de descendre à
$\mathbb{Q}$ en rendant les points de $S$ rationnels : les envoyer sur $0, 1,
\infty$. On est conduit à la \emph{tour de Fermat} :
\[
\mathcal{G}_{n} = (\Gamma \times \mathfrak{S}_{S}) \ltimes \pi_{n},
\qquad \pi_{n} = \mu_{n}^{S}/\mu_{n} .
\]
C'est le « cas favorable » de la page 7, et un cas où l'extension est
scindée ; la section 3 le développe.

\textbf{D)} Si $\operatorname{Card} S \geqslant 4$, un triplet $s = (s_{0},
s_{1}, s_{\infty})$ de points distincts de $S$ définit un isomorphisme
$\varphi_{s} : \hat{\mathfrak{X}} \simeq \mathbb{P}^{1}_{\bar{\mathbb{Q}}}$ qui
les envoie sur $0, 1, \infty$ ; les autres points deviennent des éléments de
$\bar{\mathbb{Q}} \smallsetminus \{0, 1\}$, qui engendrent un corps de nombres
$K_{s}$. Ce corps ne dépend pas de $s$ (il est engendré par les birapports des
points de $S$) : c'est le plus petit sous-corps $K$ de $\bar{\mathbb{Q}}$ sur
lequel $\hat{\mathfrak{X}}$ descend en $\hat{X}_{K}$ de façon que les points
de $S$ deviennent rationnels, et cette descente est unique à isomorphisme
unique près. Les automorphismes de $(\hat{\mathfrak{X}}, S)$ forment un groupe
fini, qui s'injecte dans $\mathfrak{S}_{S}$ et opère déjà sur
$\hat{X}_{K}$.\footnote{Ce dernier point : un automorphisme de
$(\hat{\mathfrak{X}}, S)$ est l'homographie qui envoie un triplet de $S$ sur un
autre triplet de $S$, et, dans la coordonnée $\varphi_{s}$, ses coefficients
sont dans $K$. La page s'arrête là.}

\emph{Tours à étudier} (page 17, feuillet isolé) : la tour des courbes de
Fermat $C_{n}$ ; la tour des revêtements abéliens de l'octaèdre, vu comme
carte régulière triangulaire ; les revêtements abéliens de $\mathbb{P}^{1}
\smallsetminus \{0, 1, \infty\}$, diédraux et « paradiédraux ».\footnote{« 3-gonale
(2) » est lu avec un doute ; le sens de « paradiédraux » n'est pas précisé.}

\pagerange{18}{22}
\section*{3. Calculs : $z \mapsto z^{n}$, la fonction $g$, Kummer sur $\mathbb{Q}$
(pages 18 à 22)}

Ces pages sont des feuillets de calcul ; on en dégage ce qui sert ensuite.

\emph{Les points spéciaux.} Le groupe $\mathfrak{S}_{3}$ des homographies qui
permutent $\{0, 1, \infty\}$ a, outre ces trois points, des points
remarquables : les deux points fixes des 3-cycles, $-j$ et $-\bar{j}$ (racines
de $z^{2} - z + 1$), et les points fixes non triviaux des trois
transpositions, $2$, $-1$ et $\tfrac12$. Sur la sphère de Riemann, ces sept
points sont les sommets, les milieux de côtés et les centres des deux
triangles (demi-plans supérieur et inférieur) de la triangulation de la sphère
par la droite réelle --- la figure qui reviendra comme « barycentrique ».
La fonction rationnelle de degré $6$
\[
g(z) = \frac{27\, z^{2}(z-1)^{2}}{4\,(z^{2} - z + 1)^{3}}
\]
est $\mathfrak{S}_{3}$-invariante, c'est le quotient de $\mathbb{P}^{1}$ par
$\mathfrak{S}_{3}$ : elle envoie $0, 1, \infty$ sur $0$ (ramification $2$),
$\tfrac12, 2, -1$ sur $1$ (ramification $2$), $-j, -\bar{j}$ sur $\infty$
(ramification $3$), grâce à l'identité
\[
4(z^{2} - z + 1)^{3} - 27 z^{2}(z-1)^{2} = \bigl[(z+1)(2z-1)(z-2)\bigr]^{2},
\]
qui est sur la page et que l'on vérifie directement.\footnote{Le nom de
« fonction de Belyi » est le nôtre ; le théorème de Belyi est de 1979. Au
facteur près, $g$ est l'inverse de l'invariant modulaire exprimé en le
paramètre de Legendre, $j(\lambda) = 256\,(\lambda^{2} - \lambda +
1)^{3}/\lambda^{2}(\lambda-1)^{2}$.} La page la compose avec $z \mapsto z^{n}$,
$g_{n}(z) = g(z^{n})$.

\emph{Le revêtement de Kummer $z \mapsto z^{n}$.} C'est un revêtement principal
de $\mathbb{P}^{1}_{\mathbb{Q}}$ de groupe $\mu_{n}$, et le groupe de ses
automorphismes semi-linéaires est $\mathcal{E}_{n} = \Gamma \ltimes
\mu_{n}(\bar{\mathbb{Q}})$. Les images réciproques des points spéciaux sont
\[
S_{0} = \{0\}, \quad S_{1} = \mu_{n}, \quad S_{\infty} = \{\infty\}, \quad
E^{+} = \{z^{n} = -j\}, \quad E^{-} = \{z^{n} = -\bar{j}\},
\]
\[
R = \{z^{n} = -1\}, \quad R' = \{z^{n} = \tfrac12\}, \quad R'' = \{z^{n} = 2\}.
\]
Le corps engendré par $E \cup R = \{z^{3n} = -1\}$ est $\mathbb{Q}(\mu_{6n})$, qui
contient $\mathbb{Q}(S_{0} \cup S_{1} \cup S_{\infty}) = \mathbb{Q}(\mu_{n})$ ;
celui qu'engendrent $R'$ et $R''$ est $\mathbb{Q}(\mu_{n}, \sqrt[n]{2})$, de degré
$\nu$ sur $K_{n} = \mathbb{Q}(\mu_{n})$, où $\nu$ est l'ordre de $2$ dans
$K_{n}^{*}/K_{n}^{*n}$.\footnote{La page écrit $\mathbb{Q}[R'] =
\mathbb{Q}[\sqrt[n]{2}]$ ; le corps engendré par \emph{toutes} les racines
$n$-ièmes de $\tfrac12$ contient leurs quotients, donc $\mu_{n}$.} Ces six
ensembles $E^{+}, E^{-}, R, R', R'', S_{1}$ sont, à la limite, des
$\mathrm{T}$-torseurs $K(-j), K(-\bar{j}), K(-1), K(\tfrac12), K(2), K(1)$, avec
les relations
\[
E^{+} \wedge E^{-} = 1, \qquad R' \wedge R'' = 1, \qquad R^{2} = 1, \qquad
(E^{+})^{6} = 1
\]
dans le groupe des classes de $\mathrm{T}$-torseurs, qui viennent de
$(-j)(-\bar{j}) = 1$, $2 \cdot \tfrac12 = 1$, $(-1)^{2} = 1$, $(-j)^{6} = 1$ ;
l'inversion $z \mapsto 1/z$ échange $E^{+}$ et $E^{-}$, $R'$ et $R''$, $S_{0}$ et
$S_{\infty}$.\footnote{Page 22 ; la page écrit $S_{2}$ pour $S_{\infty}$.}

\emph{Kummer sur $\mathbb{Q}$.} La théorie de Kummer donne
\[
\mathbb{Q}^{*}/\mathbb{Q}^{*n} \xrightarrow{\;\sim\;} H^{1}(\Gamma, \mu_{n}),
\]
$\xi$ s'envoyant sur la classe du cocycle $\gamma \mapsto
\gamma(\sqrt[n]{\xi})/\sqrt[n]{\xi}$, et $\sqrt[n]{\xi}$ est rationnel si et
seulement si ce cocycle est nul. Les cocycles eux-mêmes, $Z^{1}(\Gamma,
\mu_{n})$, forment une extension de $H^{1}$ par les cobords $B^{1} \simeq
\mu_{n}/\mu_{n}(\mathbb{Q})$ : choisir la racine, c'est choisir le cocycle dans
sa classe. Passant à la limite, $\varprojlim_{n} \mathbb{Q}^{*}/\mathbb{Q}^{*n}
\simeq \{\pm 1\} \times \varprojlim_{n} (\mathbb{Z}/n)^{(P)}$, $P$ l'ensemble des
nombres premiers, et le second facteur s'injecte dans
$\widehat{\mathbb{Z}}^{P}$.\footnote{La page écrit les coefficients
$\mathbb{Z}/n\mathbb{Z}$ et $\widehat{\mathbb{Z}}$ (pp.~20--21) ; ils ne valent
qu'après l'identification $\mu_{n} \simeq \mathbb{Z}/n$ par
$\exp(2i\pi/n)$, qui n'est pas $\Gamma$-équivariante. Elle écrit aussi
$Z^{1}(\Gamma, \mathbb{Z}/n) \simeq \mathbb{Q}^{*}/\mathbb{Q}^{*n} \oplus
\mathbb{Z}/n$ ; le noyau de $Z^{1} \to H^{1}$ est $\mu_{n}/\mu_{n}(\mathbb{Q})$,
qui n'est $\mathbb{Z}/n$ que pour $n$ impair, comme le dit une ligne de la page
20. La flèche $\varprojlim (\mathbb{Z}/n)^{(P)} \to \widehat{\mathbb{Z}}^{P}$ est
injective, non surjective.} D'où, pour $\xi = 2$, un homomorphisme
\[
\Gamma \longrightarrow \mathrm{Aff}(1, \widehat{\mathbb{Z}}) =
\widehat{\mathbb{Z}}^{*} \ltimes \widehat{\mathbb{Z}},
\]
action de $\Gamma$ sur le torseur $K(2) \simeq \widehat{\mathbb{Z}}$, et, pour
plusieurs rationnels, $\Gamma \to \widehat{\mathbb{Z}}^{*} \ltimes
\widehat{\mathbb{Z}}^{n}$. La page demande si l'image est d'indice fini : elle
l'est pour des rationnels multiplicativement indépendants, et pour $\xi = 2$
elle est d'indice $2$ (voir la page 35).\footnote{Réponse nôtre : c'est un
résultat classique de la théorie de Kummer sur $\mathbb{Q}$ : le groupe de Galois de
$\mathbb{Q}(\mu_{\infty}, \xi_{1}^{1/\infty}, \ldots, \xi_{n}^{1/\infty})$ sur
$\mathbb{Q}(\mu_{\infty})$ est un sous-groupe ouvert de $\widehat{\mathbb{Z}}^{n}$. La page 19 pose aussi une
question sur les sous-tores finis $D \subset \mathbb{G}_{m}$ et
$\underline{\mathrm{Hom}}(D, \mathbb{G}_{m})$, que l'on ne sait pas
reconstituer : plusieurs mots en sont illisibles.}

La page 22 esquisse le même calcul pour le revêtement de Kummer modifié par
les points $R', R''$ : une extension $1 \to \mu_{n} \to \mathcal{E}'_{n} \to
\Gamma' \to 1$, où $\Gamma'$ est un quotient de $(\mathbb{Z}/6n)^{*} \ltimes
\mu_{\nu}$, scindée, le scindage dépendant du choix de $\sqrt[\nu]{2}$.%
\footnote{Lecture d'ensemble incertaine : l'indice de $\Gamma'$ en bas du
diagramme est illisible.}

\section*{4. La tour de Fermat (pages 24 à 40)}

\pagerange{24}{28}
\subsection*{Le revêtement abélien universel de $\mathbb{P}^{1} \smallsetminus
\{0, 1, \infty\}$ (§§ 1.1 à 1.5, pages 24 à 28)}

Soit $C_{n} \subset \mathbb{P}^{2}_{\mathbb{Q}}$ la courbe de Fermat $x^{n} +
y^{n} + z^{n} = 0$. Le groupe
\[
\Pi_{n} = \mu_{n}^{3}/\mu_{n} \simeq \mu_{n} \otimes_{\mathbb{Z}}
(\mathbb{Z}^{3}/\mathbb{Z})
\]
(quotient par la diagonale) opère sur $C_{n}$ par matrices diagonales, et
$(x : y : z) \mapsto (x^{n} : y^{n} : z^{n})$ identifie $C_{n}/\Pi_{n}$ à la droite
$C_{1} : X + Y + Z = 0$. On identifie $C_{1}$ à $\mathbb{P}^{1}$ par la
coordonnée $t = -X/Z$, qui envoie les points où $X$, $Y$, $Z$ s'annulent sur
$0$, $1$, $\infty$. Le revêtement $C_{n} \to C_{1}$ est étale hors de ces trois
points, où il a $n$ points d'indice de ramification $n$ : c'est un revêtement
principal $C_{n}^{*} \to C_{1}^{*} = \mathbb{P}^{1} \smallsetminus \{0, 1,
\infty\}$ de groupe $\Pi_{n}$. Topologiquement, les deux demi-plans de $C_{1}$
se relèvent en une triangulation de $C_{n}$ par $2n^{2}$ triangles. À la limite,
\[
C_{\infty} = \varprojlim C_{n}, \qquad
\Pi_{\infty} = \varprojlim \Pi_{n} \simeq \mathrm{T}^{3}/\mathrm{T} \simeq
\mathrm{T} \otimes_{\widehat{\mathbb{Z}}} (\widehat{\mathbb{Z}}^{3}/\widehat{\mathbb{Z}}).
\]

Le groupe $\mathfrak{S}_{3}$ permute les coordonnées ; il opère sur chaque
$C_{n}$, sur $\Pi_{n}$, compatiblement, et sur $C_{1}$ par son action sur $\{0,
1, \infty\}$. Après extension à $\bar{\mathbb{Q}}$, $\Gamma$ opère sur les
$\bar{C}_{n}$ et sur $\bar{\Pi}_{n} = \Pi_{n}(\bar{\mathbb{Q}})$ par le caractère
$\chi_{n}$, et ces actions se recollent en
\[
\mathcal{E}_{n} = (\mathfrak{S}_{3} \times \Gamma) \ltimes \bar{\Pi}_{n},
\qquad
\mathcal{E}_{\infty} = (\mathfrak{S}_{3} \times \Gamma) \ltimes \Pi_{\infty},
\]
où $\mathfrak{S}_{3}$ opère sur $\mathbb{Z}^{3}/\mathbb{Z}$ par permutation et
$\Gamma$ par l'homothétie $\chi$. On a
\[
\mathfrak{S}_{3} \times \Gamma = \mathrm{Aut}_{\mathbb{Q}}(\bar{C}_{1}, \{0, 1,
\infty\}), \qquad
\mathcal{E}_{n} = \mathrm{Aut}_{\mathbb{Q}}(\bar{C}_{n} \to \bar{C}_{1}),
\]
et $\mathcal{E}_{n}$ est le groupe de tous les $\mathbb{Q}$-automorphismes de
$\bar{C}_{n}$ dès que $n \geqslant 4$.\footnote{La page (13) écrit
$\mathfrak{S}_{3} \times \Gamma = \mathrm{Aut}_{\mathbb{Q}}\, C_{1}$ et
$\mathcal{E}_{n} \simeq \mathrm{Aut}_{\mathbb{Q}}(C_{n})$, avec $\mu_{n}$ là où
(9) a $\Pi_{n}$. Pour $n \leqslant 3$, $C_{n}$ est de genre $0$ ou $1$ et a
d'autres automorphismes ; pour $n \geqslant 4$, le groupe des automorphismes
géométriques de la courbe de Fermat est $\mathfrak{S}_{3} \ltimes \Pi_{n}$
d'ordre $6n^{2}$. Les groupes $\mathcal{E}_{n}$ sont le $\mathcal{G}'$ de la
page 6 pour le couple $\bar{C}_{n} \to \bar{C}_{1}$, avec $\mathfrak{G} =
\mathfrak{S}_{3}$ et $G = \bar{\Pi}_{n}$.} Enfin l'homomorphisme
$\pi_{1}(\mathbb{P}^{1}_{\bar{\mathbb{Q}}} \smallsetminus \{0, 1, \infty\}) \to
\Pi_{\infty}$ induit un isomorphisme
\[
\pi_{1}(\mathbb{P}^{1}_{\bar{\mathbb{Q}}} \smallsetminus \{0, 1,
\infty\})^{\mathrm{ab}} \simeq \Pi_{\infty} :
\]
la tour de Fermat est le revêtement abélien universel, et « on est en train de
faire l'étude arithmétique des revêtements étales abéliens » de la droite
privée de trois points.\footnote{L'isomorphisme est $\Gamma$-équivariant, ce
que dit la forme $\mathrm{T} \otimes (\widehat{\mathbb{Z}}^{3}/\widehat{\mathbb{Z}})$ :
le $\pi_{1}^{\mathrm{ab}}$ est $H_{1} \otimes \widehat{\mathbb{Z}}$ tordu à la
Tate. La réalisation de ces courbes comme revêtements abéliens universels
sera exploitée par Y.~Ihara (« Profinite braid groups, Galois representations
and complex multiplications », 1986), qui part des jacobiennes des courbes de
Fermat ; ce travail est postérieur.}

\pagerange{28}{33}
\subsection*{Les fibres au-dessus des sommets et des centres (§§ 1.6 et 1.7,
pages 28 à 33)}

Au-dessus de $0$, la fibre de $\bar{C}_{n}$ est
$\{(0 : \lambda : 1) \mid \lambda^{n} = -1\}$, et à la limite
\[
E_{0} \simeq K(-1) = p_{2}^{-1}(-1), \qquad p_{2} : \mathrm{T} \to \mu_{2},
\]
le $\mathrm{T}$-torseur des systèmes compatibles de racines de $-1$,
classe de $\mathrm{T}$ modulo $2\mathrm{T}$. Il en va de même au-dessus de
$1$ et de $\infty$, par $\mathfrak{S}_{3}$. Pour une description intrinsèque,
compatible à $\mathfrak{S}_{3}$, $\Gamma$ et $\Pi_{\infty}$, considérons la suite
exacte de $(\mathfrak{S}_{3} \times \Gamma)$-modules
\[
0 \to \Pi_{\infty} \xrightarrow{\;2\;} \Pi_{\infty} \xrightarrow{\;p\;} \Pi_{2} \to 0,
\]
et remarquons que $\Pi_{2} \smallsetminus \{0\} = (\mathbb{F}_{2}^{3}/\mathbb{F}_{2})
\smallsetminus \{0\}$ a trois éléments, en bijection $\mathfrak{S}_{3}$-équivariante
avec $I = \{0, 1, \infty\}$, et que $\Gamma$ y opère trivialement. Pour $i \in
I$, soit $Z_{i} \subset \Pi_{\infty}$ le sous-groupe d'inertie, image du facteur
$\mathrm{T}$ d'indice $i$ de $\mathrm{T}^{3}$. Alors
\[
E_{i} \simeq p^{-1}(i)/2Z_{i},
\]
$\Pi_{\infty}$ opérant sur la classe $p^{-1}(i)$ par $y \mapsto y + 2x$ : c'est
un $\Pi_{\infty}/Z_{i}$-torseur, comme doit l'être la fibre au-dessus d'un
point de ramification d'inertie $Z_{i}$.\footnote{La page (20) écrit
$p^{-1}(\{i\})/Z_{i}$. Mais $Z_{i}$ ne stabilise pas la classe $p^{-1}(i)$ de
$\Pi_{\infty}$ modulo $2\Pi_{\infty}$ (son générateur $\rho_{i}$ a pour image $i$
dans $\Pi_{2}$), et ce sont les translations par $2\Pi_{\infty}$ qui opèrent ;
le quotient juste est par $2Z_{i} = Z_{i} \cap 2\Pi_{\infty}$. Une note de marge
propose une autre voie, par les trois points rationnels.}

Soient maintenant $\alpha = -j$ et $\bar{\alpha} = -\bar{j}$, les deux points de
$\bar{C}_{1}$ fixes sous $\mathfrak{S}_{3}^{+}$, et $\alpha_{*} = \{\alpha,
\bar{\alpha}\}$. Pour décrire $E(\alpha_{*})$, la fibre au-dessus de
$\alpha_{*}$, avec ses opérations, on passe aux cônes : soit
\[
\widetilde{C}_{1} = \mathrm{Isom}(I, \mu_{3}) \simeq \{(\zeta_{0}, \zeta_{1},
\zeta_{\infty}) \in \mu_{3}^{3} \mid \text{les } \zeta_{i} \text{ distincts}\},
\]
l'ensemble à $6$ éléments des points du cône de $\bar{C}_{1}$ à coordonnées
racines cubiques de l'unité distinctes (leur somme est nulle). Le groupe
$\mu_{3}$ y opère librement par homothéties, et $\widetilde{C}_{1}/\mu_{3}
\xrightarrow{\sim} \alpha_{*}$ : deux éléments, les deux ordres cycliques de
$I$. Le groupe $\mathcal{E}_{\infty}$ y opère à travers $\mathfrak{S}_{3}
\times \{\pm 1\}$, via $\chi_{3}$. Soit $\widetilde{C}_{n}$ l'image réciproque de
$\widetilde{C}_{1}$ dans le cône de $\bar{C}_{n}$,
\[
\widetilde{C}_{n} = \{(x, y, z) \mid x^{3n} = y^{3n} = z^{3n} = 1,\ x^{n}, y^{n},
z^{n} \text{ distincts}\} ;
\]
$\widetilde{C}_{n} \to \widetilde{C}_{1}$ est compatible à $\mu_{3n}
\xrightarrow{\zeta \mapsto \zeta^{n}} \mu_{3}$, d'où $\widetilde{C}_{n}/\mu_{3n}
\simeq E_{n}(\alpha_{*}) \to \widetilde{C}_{1}/\mu_{3} \simeq
E_{1}(\alpha_{*})$.\footnote{La page écrit $\mu_{n}$ au lieu de $\mu_{3}$ dans
$\widetilde{C}_{1}/\mu_{n}$ et dans le diagramme (27) ; $\zeta \mapsto
\zeta^{n}$ envoie $\mu_{3n}$ sur $\mu_{3}$.} À la limite, avec $\xi \mapsto
\xi_{3}$ la réduction $\mathrm{T} \to \mathrm{T}/3\mathrm{T} \simeq \mu_{3}$,
\[
\widetilde{C}_{\infty} = \{(x, y, z) \in \mathrm{T}^{3} \mid x_{3}, y_{3}, z_{3}
\text{ distincts}\}, \qquad
E(\alpha_{*}) \simeq \widetilde{C}_{\infty}/\mathrm{T} = p_{3}^{-1}(\omega)
\subset \Pi_{\infty},
\]
où $p_{3} : \Pi_{\infty} \to \Pi_{3}$ et $\omega = \widetilde{C}_{1}/\mu_{3}
\subset \Pi_{3}$. Le groupe $\Pi_{\infty}$ opère librement sur
$p_{3}^{-1}(\omega)$ par $y \mapsto y + 3x$, et $\Gamma$, $\mathfrak{S}_{3}$ par
leur action sur $\Pi_{\infty}$.\footnote{La page dit que $\Pi_{\infty}$ « opère
librement sur cette partie » sans dire comment ; les translations par
$\Pi_{\infty}$ ne la stabilisent pas, celles par $3\Pi_{\infty}$ si : c'est
l'action de $\Pi_{n}$ sur $\widetilde{C}_{3n}$ par $\mu_{n}^{3} \subset
\mu_{3n}^{3}$, passée à la limite.} Sur $\Pi_{3} \simeq \mathbb{F}_{3}^{2}$, le
sous-groupe $\mathfrak{S}_{3}^{+} \simeq \mathbb{Z}/3$ opère par un bloc de
Jordan de taille $2$ : il a une unique droite invariante, et $\omega$ en est
formé des deux éléments non nuls.\footnote{Avec la coordonnée $t = -X/Z$ de la
page 24, le point $(\bar{j} : j : 1)$ que la page attache à $\alpha$ (p.~37)
est $t = -\bar{j}$, et les formules (23) décrivent la fibre au-dessus de
$-\bar{j}$. Les deux points étant échangés par la conjugaison complexe, qui
commute à tout le reste, on garde le nom de la page.}

\pagerange{34}{36}
\subsection*{Les fibres au-dessus des milieux de côtés (§ 1.8, pages 34 à 36)}

Soient $\alpha'(0) = 2$, $\alpha'(1) = -1$, $\alpha'(\infty) = \tfrac12$ les
points fixes, autres que $0, 1, \infty$, des transpositions $\tau_{0}, \tau_{1},
\tau_{\infty}$, et $\alpha'_{*}$ leur ensemble.\footnote{La page écrit d'abord
$\beta_{0}, \beta_{1}, \beta_{\infty}$, corrigés en $\alpha'(i)$, et une fois
« $\alpha'(2)$ » pour $\alpha'(\infty)$ ; le « $\beta_{*}$ » de la page 39 est
un reste de l'ancienne lettre.} Un calcul en coordonnées affines donne
\[
E_{n}(\alpha'(i)) \simeq \{(\xi_{j})_{j \in I \smallsetminus \{i\}} \mid
\xi_{j}^{n} = -\tfrac12\}, \qquad
E(\alpha'(i)) \simeq K(-\tfrac12)^{I \smallsetminus \{i\}},
\]
et les torseurs $K(g)$ se composent : $K(gg') \simeq K(g) \wedge_{\mathrm{T}}
K(g')$, d'où $K(-\tfrac12) \simeq K(-1) \wedge_{\mathrm{T}} K(2)^{-1}$, avec
$K(-1) \subset \mathrm{T}$ comme plus haut. Toute l'arithmétique de ces fibres
est donc dans les deux torseurs $K(-1)$, sur lequel $\Gamma$ opère à travers
$\chi$, et $K(2)$.

Soit $A = \mathrm{Aff}(\mathrm{T}, K(\tfrac12))$ le groupe des automorphismes
affines du torseur $K(\tfrac12)$ ; c'est une extension
\[
1 \to \mathrm{T} \to A \to \widehat{\mathbb{Z}}^{*} \to 1 .
\]
Le groupe $\widetilde{\mathcal{E}} = (\mathfrak{S}_{3} \times A) \ltimes
\Pi_{\infty}$, où $A$ opère sur $\Pi_{\infty}$ à travers $\widehat{\mathbb{Z}}^{*}$,
opère sur toutes les fibres $E(\alpha_{*})$, $E(\{0, 1, \infty\})$,
$E(\alpha'_{*})$ ; sur les deux premières à travers $\widetilde{\mathcal{E}}/\mathrm{T}
\simeq (\mathfrak{S}_{3} \times \widehat{\mathbb{Z}}^{*}) \ltimes \Pi_{\infty}$.
L'action de $\Gamma$ sur $K(\tfrac12)$ donne $\Gamma \to A$, dont l'image
$\Gamma'$ est le groupe de Galois de $\mathbb{Q}(\mu_{\infty}, 2^{1/\infty})$, et
$\mathcal{E}_{\infty} \to \widetilde{\mathcal{E}}$, à travers lequel se font
toutes les opérations précédentes. La page demande si $\Gamma \to A$ est
surjectif : il ne l'est pas, son image est d'indice $2$.\footnote{Réponse
nôtre. Comme $\sqrt{2} = \zeta_{8} + \zeta_{8}^{-1} \in \mathbb{Q}(\mu_{8})$, la
racine carrée de $2$ est déterminée par l'action sur les racines de l'unité ;
c'est la seule obstruction, et l'image est d'indice exactement $2$.}

\pagerange{36}{40}
\subsection*{Trivialisation en $\alpha$, et le transport parallèle (§ 1.9,
pages 36 à 40)}

Rien de ce qui précède n'utilise le plongement $\bar{\mathbb{Q}} \subset
\mathbb{C}$ : tout reste valable sur un corps premier quelconque, pour les $n$
premiers à la caractéristique, $\bar{C}_{\infty}$ étant alors le revêtement
abélien universel \emph{modéré}. Le plongement sert à faire le lien « avec les
cartes finies abéliennes »,\footnote{« 3-gradiées » et « point cartographique »
sont lus avec un doute.} en trivialisant $\bar{C}_{\infty}$ au-dessus de
$\alpha$. On pose
\[
j = \exp\tfrac{2i\pi}{3}, \qquad \zeta_{n} = \exp\tfrac{2i\pi}{3n}, \qquad
\alpha_{n} = (\bar{\zeta}_{n}, \zeta_{n}, 1) \in \widetilde{C}_{n}, \qquad
\alpha_{\infty} = (\bar{\zeta}_{\infty}, \zeta_{\infty}, 1),
\]
où $\zeta_{\infty} = (\zeta_{n})$ est le générateur de $\mathrm{T}$ donné par
l'exponentielle complexe, d'où $\varphi : \widehat{\mathbb{Z}} \simeq
\mathrm{T}$, et
\[
\Pi_{\infty} \simeq \widehat{\mathbb{Z}}^{3}/\widehat{\mathbb{Z}}, \qquad
\mathcal{E}_{\infty} \simeq (\mathfrak{S}_{3} \times \Gamma) \ltimes
(\widehat{\mathbb{Z}}^{3}/\widehat{\mathbb{Z}}),
\]
$\mathfrak{S}_{3}$ permutant les facteurs et $\Gamma$ opérant par homothétie de
rapport $\chi$. Dans ces coordonnées $\alpha_{\infty} = (-1, 1, 0)$ et
$\bar{\alpha}_{\infty} = (1, -1, 0)$ modulo $\widehat{\mathbb{Z}}$, et les
générateurs d'inertie
\[
\rho_{0} = (1,0,0), \quad \rho_{1} = (0,1,0), \quad \rho_{\infty} = (0,0,1)
\pmod{\widehat{\mathbb{Z}}}, \qquad \rho_{0} + \rho_{1} + \rho_{\infty} = 0,
\]
« donnent l'isomorphisme » de $\Pi_{\infty}$ avec l'abélianisé du groupe
profini engendré par trois lacets de produit $1$.\footnote{La note de marge
porte « $\widehat{\Gamma}^{+}_{12}$ », indice lu sans assurance. La page
(47') appelle $\alpha'_{\infty}$ ce que l'on note $\bar{\alpha}_{\infty}$ :
l'$\alpha'$ du § 1.8 est autre chose.} En choisissant l'origine $\theta_{\infty}
= (\sqrt[n]{2})_{n}$ (racines réelles positives), $K(2) \simeq \widehat{\mathbb{Z}}$
et $A \simeq \mathrm{Aff}(1, \widehat{\mathbb{Z}}) = \widehat{\mathbb{Z}}^{*}
\ltimes \widehat{\mathbb{Z}}$ opère par $x \mapsto ax + b$ sur $K(2)$ et par $x
\mapsto ax - b$ sur $K(\tfrac12) = K(2)^{-1}$ ; tout est alors explicite.

Reste le \emph{transport parallèle} : sur $\mathbb{C}$, un chemin du demi-plan
qui contient $\alpha$ vers chacun des six autres points spéciaux identifie la
fibre $E(\alpha)$ aux fibres $E_{0}, E_{1}, E_{\infty}$ et $E'_{0} = E_{2}$,
$E'_{1} = E_{-1}$, $E'_{\infty} = E_{1/2}$ --- dans le langage des cartes, il dit
à quels sommets et à quels côtés touche chaque triangle. Ces bijections étant
compatibles à $\Pi_{\infty}$ et à $\mathfrak{S}_{3} \times \{1, \tau\}$ ($\tau$ la
conjugaison complexe), il suffit d'en connaître deux, par exemple $E(\alpha) \to
E'_{\infty} \to E_{0}$, c'est-à-dire l'image du point marqué $\alpha_{\infty}$.
La page calcule le point de $\bar{C}_{n}$ au-dessus de $\tfrac12$,
\[
\alpha'_{n}(\infty) = \Bigl(\sqrt[n]{\tfrac12}\, e^{-i\pi/n},\
\sqrt[n]{\tfrac12}\, e^{i\pi/n},\ 1\Bigr),
\]
et s'arrête.\footnote{La ligne suivante de la page récrit ce point
$(\theta_{n}^{-1}\bar{\zeta}_{n}, \theta_{n}^{-1}\zeta_{n}, 1)$ ; avec
$\zeta_{n} = \exp(2i\pi/3n)$ de (43), $e^{i\pi/n}$ n'est pas $\zeta_{n}$, et
l'on ne garde que la première ligne, qui est juste ($(\sqrt[n]{1/2}\,
e^{-i\pi/n})^{n} = -\tfrac12$). L'argument s'interrompt au bas de la page 40.}

\section*{5. Fermat, suite : trois points abstraits, l'octaèdre, la tour diédrale
(pages 41 à 54)}

\pagerange{41}{41}
\subsection*{Fonctorialité en un ensemble à trois éléments (page 41)}

La page 41 reprend ce qui précède en remplaçant $\{0, 1, \infty\}$ par un
ensemble $I$ à trois éléments quelconque : $\pi(I) =
\widehat{\mathbb{Z}}^{I}/\widehat{\mathbb{Z}}$, $\pi_{n}(I) =
(\mathbb{Z}/n)^{I}/(\mathbb{Z}/n)$, le groupe $(\mathfrak{S}_{I} \times
\widehat{\mathbb{Z}}^{\times}) \ltimes \pi(I)$, les fibres $E(I, i)$ au-dessus
des points de $I$ et leur réunion $E(I; I) = \coprod_{i} E(I, i)$, et la fibre
$E(I, \omega) = p_{3}^{-1}(\omega) \subset \pi_{\infty}(I)$ au-dessus des deux
« centres », où $\omega \simeq \mathrm{Or}(I)$ est l'ensemble des deux ordres
cycliques de $I$. C'est la forme intrinsèque de la page 33 : l'ensemble
$\omega$ des deux points fixes des 3-cycles est l'ensemble des deux
orientations du triangle.\footnote{Plusieurs lectures de la page sont
douteuses ($\bar{\jmath}$ dans l'accolade, le premier terme de la suite
exacte).}

\pagerange{43}{43}
\subsection*{L'octaèdre (page 43)}

Le genre de $C_{n}$, courbe plane lisse de degré $n$, est $(n-1)(n-2)/2$ : nul
pour $n = 1, 2$, égal à $1$ pour $n = 3$, à $3$, $10$, $55$ pour $n = 4, 6, 12$.
Pour $n = 2$, la conique $C_{2} : x^{2} + y^{2} + z^{2} = 0$, avec sa
triangulation, est formée de $8$ triangles, chaque sommet en touchant $4$ :
c'est l'\emph{octaèdre}, dont le groupe des rotations est $\mathfrak{S}_{4}
\simeq \mathfrak{S}_{3} \ltimes (\mathbb{Z}/2)^{2}$, c'est-à-dire
$\mathfrak{S}_{3} \ltimes \Pi_{2}$. Les quotients de $C_{2}$ par les trois
sous-groupes d'ordre $2$ de $\Pi_{2}$ sont les trois revêtements doubles de
$C_{1}$ ramifiés en deux des trois points, permutés par $\mathfrak{S}_{3}$.
L'octaèdre a le type de ramification $(2, 2), (2, 2), (2, 2)$ au-dessus de $0, 1,
\infty$, qui est celui du revêtement diédral de type $(2, 2, n)$ pour $n = 2$ ;
comme un revêtement galoisien de groupe $(\mathbb{Z}/2)^{2}$ ramifié en trois
points avec toutes ses inerties non triviales est unique, les deux
coïncident.\footnote{La note encadrée de la page se lit mal ; on en restitue le
sens, sans ses mots (« subdivision barycentrique du biédre », « autodual »,
« unicité de la rev.-théorie » sont incertains).}

\pagerange{45}{46}
\subsection*{La tour diédrale (pages 45 et 46)}

Sur $X = C_{2}$, le diviseur de ramification $D$ est de degré $6$ et défini sur
$\mathbb{Q}(i)$. Les pages 45--46 construisent une seconde tour, en coordonnées :
\[
X_{n} \xrightarrow{\;x_{n} \mapsto x_{n}^{n}\;} X_{1} \xrightarrow{\;z \mapsto
x_{0}\;} X_{0}, \qquad
x_{0} = \tfrac12 - \tfrac14\bigl(z + z^{-1}\bigr).
\]
Le revêtement $X_{1} \to X_{0}$ est de degré $2$, ramifié en $z = 1$ (au-dessus
de $x_{0} = 0$) et en $z = -1$ (au-dessus de $x_{0} = 1$), non ramifié au-dessus
de $\infty$ ($z = 0, \infty$) ; $X_{n} \to X_{1}$ est ramifié d'indice $n$ en $0$ et
$\infty$. Le composé $X_{n} \to X_{0}$, $x_{0} = \tfrac12 - \tfrac14(x_{n}^{n} +
x_{n}^{-n})$, est galoisien de groupe diédral $D_{n} = (\mathbb{Z}/2) \ltimes \mu_{n}$
et de type de ramification $(2, 2, n)$ ; à la limite, de groupe
$\mathbb{D}_{\infty} = (\mathbb{Z}/2) \ltimes \mu_{\infty}$.\footnote{Le nom de
cette tour, et le groupe $D_{n}$, sont tirés des formules de la page 46 ; le
lien avec $C_{2}$ annoncé à la page 45 n'est pas écrit, la page s'arrêtant
après « il a comme groupe $(\pi_{\infty})_{2}$ ».}

\pagerange{48}{54}
\subsection*{Les endomorphismes de la conique (pages 48 à 54)}

Sur la conique $C_{2}$, la fonction $u = (y - ix)/z$ est une coordonnée définie
sur $\mathbb{Q}(i)$, et l'élévation à la puissance $n$ dans cette coordonnée est
un endomorphisme de degré $n$. Pour $n = 3$, la page 54 calcule
\[
\Bigl(\frac{y - ix}{z}\Bigr)^{3} = \frac{1}{z^{3}}\Bigl[\,y(y^{2} - 3x^{2}) +
i\, x(x^{2} - 3y^{2})\,\Bigr],
\]
ce qui retrouve l'application de la page 48,
\[
(x, y, z) \longmapsto \bigl(-x(x^{2} - 3y^{2}),\ y(y^{2} - 3x^{2}),\ z^{3}\bigr),
\]
dont la page vérifie qu'elle préserve la conique : la somme des carrés des
images vaut $x^{6} + y^{6} + z^{6} + 3x^{2}y^{2}(x^{2} + y^{2})$, c'est-à-dire, sur
$C_{2}$, $a^{3} + b^{3} + c^{3} - 3abc$ avec $a, b, c = x^{2}, y^{2}, z^{2}$, qui est
divisible par $a + b + c$.\footnote{La page écrit l'égalité avec $x^{6} + y^{6} +
z^{6} - 3x^{2}y^{2}z^{2}$ sans dire qu'elle n'a lieu que modulo l'équation de la
conique, et ne conclut pas que l'expression est nulle ; c'est ce qu'il fallait.
À la page 54 elle écrit $(y + ix)^{2}$ là où le calcul donne $(y - ix)^{2}$.} À la
question de la page 54 --- les $x', y', z'$ de l'endomorphisme d'ordre $n$
sont-ils des fonctions rationnelles de $x, y, z$ ? --- cette coordonnée répond
oui. La page 50 fait le même calcul par la paramétrisation rationnelle
$(1 - t^{2}, 2t, 1 + t^{2})$ et l'élévation au carré $(x^{2} - y^{2}, 2xy, x^{2} +
y^{2})$, qui préservent la forme $x^{2} + y^{2} - z^{2}$, équivalente à celle de
$C_{2}$ sur $\mathbb{Q}(i)$ par $z \mapsto iz$.\footnote{La page écrit $C_{2} :
x^{2} + y^{2} + z^{2} = 0$ et vérifie $x'^{2} + y'^{2} - z'^{2} = 0$ ; elle a
aussi « $(t + i)/(t - 1)$ » où la ligne précédente a $(t+i)/(t-i)$.} Les
pages 52--53 dessinent, à partir de là, un diagramme de revêtements $X_{1}
\leftarrow X_{2} \leftarrow X_{4} \leftarrow X_{8}$, $\widetilde{X}_{2} =
X_{2}(2, 2, 2)$, de groupes des quotients de $\mu_{n}^{3}/\mu_{n}$ ; il est trop
surchargé pour être restitué flèche par flèche.\footnote{Voir la transcription,
p.~52, qui en décrit les étiquettes sans choisir entre les flèches biffées.
La page 48 contient encore un calcul de valeurs critiques
($p - \lambda q = p' - \lambda q' = 0$).}

\section*{6. Automorphismes de groupoïdes (pages 56 à 62)}

\pagerange{57}{60}
\subsection*{Foncteurs d'un groupoïde au-dessus d'une application des objets
(pages 57 à 60)}

Soient $C$ un groupoïde, $C'$ une catégorie, $C_{0}, C'_{0}$ leurs ensembles
d'objets et $f_{0} : C_{0} \to C'_{0}$ une application. On cherche les foncteurs
$f : C \to C'$ qui l'induisent.

\textbf{Proposition.} a) \emph{Il existe un $f$ au-dessus de $f_{0}$ si et
seulement si $f_{0}$ envoie deux objets isomorphes (dans une même composante
connexe de $C$) sur des objets isomorphes.}

b) \emph{Choisissons dans chaque composante connexe $C(i)$ un objet $a_{i}$, et
pour chaque $x \in C(i)_{0}$ un isomorphisme $u_{x} : a_{i} \to x$, avec
$u_{a_{i}} = \mathrm{id}$. Alors $f \mapsto ((f(u_{x}))_{x}, (f_{a_{i}, a_{i}})_{i})$
est une bijection}
\[
\mathrm{Hom}_{f_{0}}(C, C') \xrightarrow{\;\sim\;}
\prod_{x \in C_{0},\, x \neq a_{i(x)}} \mathrm{Isom}\bigl(f_{0}(a_{i(x)}), f_{0}(x)\bigr)
\times \prod_{i} \mathrm{Hom}\bigl(\mathrm{Aut}(a_{i}), \mathrm{Aut}(f_{0}(a_{i}))\bigr).
\]
\footnote{L'indice du premier produit est surchargé et illisible sur la page.
Pour $x = a_{i}$, $f(u_{a_{i}}) = f(\mathrm{id}) = \mathrm{id}$ est imposé : le
produit porte sur les $x$ distincts des $a_{i}$, ce que confirme la formule
$\mathrm{Aut}^{\circ}(C, a) \simeq \prod_{x \neq a} \pi_{x}$ de la page 61.}
Le foncteur associé à $(\ell_{x}, \varphi_{i})$ est $g \mapsto \ell_{y}\,
\varphi_{i}(u_{y}^{-1} g u_{x})\, \ell_{x}^{-1}$ pour $g : x \to y$, et a) en
résulte.

c) \emph{$f$ est fidèle si et seulement si chaque $\varphi_{i}$ est injectif ; il
est pleinement fidèle si et seulement si} 1°) \emph{pour $i \neq j$,
$\mathrm{Hom}(f_{0}(a_{i}), f_{0}(a_{j})) = \emptyset$ ;} 2°) \emph{pour $x, y$
dans une même composante, toute flèche $f_{0}(x) \to f_{0}(y)$ est un
isomorphisme ;} 3°) \emph{chaque $\varphi_{i}$ est un isomorphisme. C'est une
équivalence si de plus} 4°) \emph{tout objet de $C'$ est isomorphe à un
$f_{0}(a_{i})$, et un isomorphisme si 1°)--3°) et} 4'°) \emph{$f_{0}$ est
bijective.}

Lorsque $C'$ est aussi un groupoïde, 2°) est automatique et 1°) dit que $f$ est
injectif sur les composantes, d'où :

\textbf{Corollaires.} \emph{Si $C'$ est un groupoïde, $f$ est un isomorphisme si
et seulement si a) $\pi_{0}(f)$ est bijectif, b) $f_{0}$ est bijective, c) chaque
$f_{a_{i}, a_{i}}$ est un isomorphisme ; c'est une équivalence si et seulement
si a) et c). Une application $f_{0}$ provient d'un isomorphisme $C \simeq C'$
si et seulement si elle est bijective, respecte dans les deux sens la relation
« être dans la même composante », et si $\mathrm{Aut}\, f_{0}(a_{i}) \simeq
\mathrm{Aut}(a_{i})$ pour tout $i$.}\footnote{La page énonce ces corollaires
pour $C'$ catégorie ; ils supposent $C'$ groupoïde, sans quoi 2°) peut
manquer.}

Soit $C$ \emph{connexe}, $a$ un objet et $\pi_{a} = \mathrm{Aut}_{C}(a)$. Toute
permutation de $C_{0}$ provient d'un automorphisme de $C$, d'où des suites
exactes
\[
1 \to \mathrm{Aut}^{\circ}(C) \to \mathrm{Aut}(C) \to \mathfrak{S}_{C_{0}} \to 1,
\qquad
1 \to \mathrm{Aut}^{\circ}(C, a) \to \mathrm{Aut}^{\circ}(C) \to \mathrm{Aut}(\pi_{a})
\to 1,
\]
où $\mathrm{Aut}^{\circ}(C)$ est formé des automorphismes qui sont l'identité sur
les objets, et $\mathrm{Aut}^{\circ}(C, a)$ de ceux qui sont de plus l'identité
sur $\pi_{a}$ ; et l'on trouve (page 61)
\[
\mathrm{Aut}^{\circ}(C, a) \simeq \prod_{x \in C_{0} \smallsetminus \{a\}} \pi_{x} .
\]
\footnote{La connexité est dans les lignes biffées du haut de la page 59 ; la
suite non biffée et la page 60 (« $\mathrm{Aut}(C) \to \mathfrak{S}_{C_{0}}$
surjective ») ne la répètent pas. Sans elle, l'image est formée des
permutations qui respectent les composantes et les types d'isomorphie des
groupes d'automorphismes, comme le dit le second corollaire. Le
$\pi_{a}$-torseur à droite $T_{x} = \mathrm{Isom}(a, x)$ et son commutant
$\mathrm{Aut}(x)$, introduits à la page 59, servent à cette identification.}

\pagerange{61}{62}
\subsection*{Le graphe $K(I)$ (pages 61 et 62)}

À quoi sert ce calcul apparaît aussitôt. Soit $I$ un ensemble muni d'une
structure de polygone d'ordre $n \geqslant 2$, $D \simeq D_{n}$ son groupe
d'automorphismes. Soit $K(I)$ le graphe à deux sommets $o, o'$ et dont les
arêtes $(\ell_{i})_{i \in I}$ joignent toutes $o$ à $o'$ --- la « carte de
l'orange » à $n$ fuseaux. $D$ opère sur $K(I)$, et une involution qui est
l'identité sur $I$ et échange $o, o'$ commute à $D$. Soit $C(I)$ le groupoïde
fondamental de $K(I)$ en ses deux sommets ; $D$ y opère par automorphismes,
transitivement sur les objets. On va regarder le normalisateur de $D$ dans
$\mathrm{Aut}(C(I))$, qui contient un élément canonique $\tau$ --- et la page 62
s'arrête là.\footnote{« Groupe fondamental » sur la page ; puisque $C_{0} = \{o,
o'\}$, c'est le groupoïde fondamental. Le cas $n = 2$ demande une donnée
supplémentaire, dont la description entre crochets est en partie illisible.}

\pagerange{63}{63}
\section*{7. La double pyramide (page 63)}

Une figure, sans texte, donne le modèle combinatoire d'une carte par ses
triangles. Les triangles de la subdivision barycentrique se partagent en
positifs $E^{+}$ et négatifs $E^{-}$ ; chaque triangle a un côté de chacun des
trois types $D_{01}, D_{12}, D_{20}$ (sommet--milieu d'arête, milieu--centre,
centre--sommet), et chaque côté borde un triangle de chaque signe : d'où des
bijections $p^{\pm}_{ij} : E^{\pm} \to D_{ij}$. Les permutations
\[
\rho_{0}^{+} = (p_{01}^{+})^{-1}\, p_{01}^{-}\, (p_{20}^{-})^{-1}\, p_{20}^{+},
\quad
\rho_{1}^{+} = (p_{12}^{+})^{-1}\, p_{12}^{-}\, (p_{01}^{-})^{-1}\, p_{01}^{+},
\quad
\rho_{2}^{+} = (p_{20}^{+})^{-1}\, p_{20}^{-}\, (p_{12}^{-})^{-1}\, p_{12}^{+}
\]
de $E^{+}$ --- franchir un côté puis un autre, c'est tourner autour du sommet
qui leur est commun --- satisfont
\[
\rho_{2}^{+}\rho_{1}^{+}\rho_{0}^{+} = \mathrm{id}_{E^{+}},
\]
identité qui se vérifie par simplification des facteurs adjacents, et
l'ensemble $S_{i}$ des sommets de type $i$ s'identifie aux orbites :
$S_{i} \simeq E^{+}/\rho_{i}^{+} \simeq E^{-}/\rho_{i}^{-}$. C'est la
description d'une carte orientée par trois permutations de produit
$1$.\footnote{Ce modèle, avec deux permutations, est celui de G.~Jones et
D.~Singerman, « Theory of maps on orientable surfaces » (\emph{Proc.\ London
Math.\ Soc.}, 1978), contemporain de ces pages ; Grothendieck le reprendra dans
l'\emph{Esquisse d'un programme} sous le nom de groupe cartographique. Les
indices du premier facteur de $\rho_{0}^{+}$ sont surchargés et rétablis sur le
modèle des deux autres.} La même figure revient, redessinée, à la page 77.

\pagerange{65}{67}
\section*{8. Sous-groupes finis de $PGL_{2}$, cas modéré (pages 64 à 67)}

Soit $X' \to X$ un revêtement galoisien de courbes, de degré $N$, ramifié
modérément au-dessus de points $i \in I$ d'indices $d_{i} \geqslant 2$. La formule
de Riemann--Hurwitz donne
\[
2 - 2g' = N \Bigl[(2 - 2g) - \sum_{i \in I} \Bigl(1 - \frac{1}{d_{i}}\Bigr)\Bigr] .
\]
Si $g' = 0$, alors $g = 0$ et
\[
\frac{2}{N} = 2 - \sum_{i} \Bigl(1 - \frac{1}{d_{i}}\Bigr) > 0 .
\]
Les solutions sont : $\operatorname{Card} I = 0$, $N = 1$ ; $\operatorname{Card}
I = 1$ est impossible ; $\operatorname{Card} I = 2$, $d_{1} = d_{2} = N$
quelconque ; $\operatorname{Card} I = 3$, $\sum 1/d_{i} > 1$, d'où les types
\[
(2, 2, p),\ N = 2p ; \qquad (2, 3, 3),\ N = 12 ; \qquad (2, 3, 4),\ N = 24 ; \qquad
(2, 3, 5),\ N = 60 ;
\]
$\operatorname{Card} I \geqslant 4$ est impossible, chaque terme $1 - 1/d_{i}$
valant au moins $\tfrac12$.\footnote{La page fait le cas $\operatorname{Card} I =
3$ en posant $p = \alpha + 2$, $q = \beta + 2$, d'où $\alpha\beta < 4$ ; et le cas
$\operatorname{Card} I \geqslant 4$ par une discussion plus longue, sur une page
67 entièrement barrée. L'argument d'une ligne est le nôtre.} C'est la
classification de Klein : groupes cycliques, diédraux, tétraédral, octaédral
et icosaédral.\footnote{Nom nôtre ; F.~Klein, \emph{Vorlesungen über das
Ikosaeder} (1884). Le titre de la page 64 est « Sous-groupes finis de
$GP(1)$ » ; on lit $PGL_{2}$.}

\emph{Le cas $(N, N)$.} En caractéristique $0$, ou si $p \nmid N$, $G \simeq
\mu_{N}(k)$ est cyclique. Si $p \mid N$, $N = p^{r}N'$, la ramification n'est
plus modérée --- on sort du cadre du titre ---, $G$ fixe un point et est contenu
dans un sous-groupe de Borel, donc produit semi-direct d'un groupe cyclique
d'ordre premier à $p$ par un $p$-groupe abélien élémentaire $(\mathbb{Z}/p)^{r}$,
sous-groupe du groupe additif.\footnote{Une note de marge écrit
$(\mathbb{Z}/p^{r}\mathbb{Z})$, exposant douteux ; un sous-groupe fini de $(k,
+)$ en caractéristique $p$ est d'exposant $p$. La page est barrée, et plusieurs
mots en sont illisibles.}

\emph{Le cas $(N, 2, 2)$}, $|G| = 2N$. En caractéristique $\neq 2$, une extension
kummérienne de degré $2$ de la base ramène à un revêtement $X'/Y'$ qui n'a plus
que deux points de ramification, donc au cas cyclique : $G$ a un sous-groupe
cyclique d'indice $2$.

\pagerange{70}{74}
\section*{9. Paradigme « gruppentheoretisch » des surfaces (pages 69 à 74)}

\emph{Point de vue topologique.} On considère des surfaces topologiques sans
bord, orientables, paracompactes, et un ensemble $\omega$ à deux éléments ; on
pose $T = \mathbb{Z} \wedge_{\mathbb{Z}/2} \omega$ (un $\mathbb{Z}$-module libre de
rang $1$ dont les deux générateurs sont les éléments de $\omega$), et une
$\omega$-orientation de $X$ est un isomorphisme du faisceau des entiers tordus
sur $X$ avec le faisceau constant $T_{X}$. On suppose chaque composante $X_{i}$
homéomorphe à $\hat{X}_{i} \smallsetminus S_{i}$, $\hat{X}_{i}$ compacte et
$S_{i}$ finie ; le couple $(\hat{X}_{i}, S_{i})$ est déterminé à isomorphisme
unique près par $X_{i}$, comme compactification par les bouts.\footnote{La
page parle du « décompactifié de Čech ». Nom nôtre : la compactification de
Freudenthal par les bouts.} La catégorie de ces surfaces
$\omega$-orientées est ainsi équivalente à celle des surfaces compactes
$\omega$-orientées munies d'une partie discrète $S$ ; de même, celle des
$(X, o, J)$, $J$ discrète, à celle des $(\hat{X}, o, S, J)$, $S$ et $J$ disjointes.

On passe aux catégories \emph{isotopiques} --- mêmes objets, morphismes les
classes d'isotopie d'homéomorphismes --- et l'on en cherche une description
par la théorie des groupes ; on se ramène au cas connexe, et l'on garde $J$
pour la dissymétrie entre points marqués et trous.

1) \emph{Cas $J = \varnothing$.} Hors du cas $g = 0$, $\nu \in \{0, 1\}$ ($g$ le
genre de $\hat{X}$, $\nu = \operatorname{card} S$), on associe à $(\hat{X},
o, S)$ le groupe fondamental $\Pi = \pi_{1}(X)$ comme groupe \emph{extérieur},
muni de sa $T$-structure : pour chaque $s \in S$, la classe du lacet autour de
$s$, c'est-à-dire une application
\[
S \longrightarrow \mathrm{Homext}(T, \Pi).
\]
Le quotient de $\Pi$ par le sous-groupe distingué engendré par ces lacets est
$\pi_{1}(\hat{X})$ ; si $\nu = 0$, la structure à donner est à la place
l'orientation $T \hookrightarrow (\bigwedge^{2} \Pi_{\mathrm{ab}})^{\vee}$. Le
foncteur ainsi obtenu, de la catégorie isotopique vers les groupes extérieurs
à $T$-structure, est une équivalence sur son image, et la description est
même « $\infty$-isotopique » --- les espaces d'homéomorphismes ont des
composantes contractiles --- sauf si $g = 0, \nu = 2$ ou $g = 1, \nu = 0$, où la
composante neutre du groupe des automorphismes est un $K(\mathbb{Z}, 1)$,
resp.\ un $K(\mathbb{Z}^{2}, 1)$.\footnote{Noms nôtres : c'est le théorème de
Dehn--Nielsen--Baer (Dehn, Nielsen 1927, Baer 1928) --- les classes d'isotopie
d'homéomorphismes d'une surface de type fini sont les automorphismes
extérieurs de son groupe fondamental qui respectent la structure
périphérique ---, et, pour la contractibilité, C.~Earle et J.~Eells (1969) et
M.-E.~Hamstrom. La page ne cite personne. Plusieurs mots de la page 73 et de
sa marge sont illisibles ; la phrase sur ce que donne la $T$-structure est
reconstituée.}

2) \emph{Cas $J \neq \varnothing$.} Une description « plus jolie » utilise le
groupoïde fondamental de $X$ en les points de $J$, avec la même $T$-structure ;
elle est $\infty$-isotopique même dans les cas $g = 0, \nu = 2$ et $g = 1, \nu =
0$.

\pagerange{76}{77}
\section*{10. Quelques problèmes spéciaux : les drapeaux (pages 75 à 77)}

La page 76 code une carte par l'ensemble $R$ de ses couples incidents $(s, f)$
(sommet, face), avec des applications $p : R \to S$, $r : R \to F$, $q', q'' : R
\to A$ (les deux arêtes du coin $(s, f)$) et des involutions $\sigma_{0}$,
$\sigma_{2}$ qui changent de face, resp.\ de sommet, autour d'une arête. Le
produit fibré
\[
E = (R, q') \times_{A} (R, q'') = \bigl\{\bigl((s, f), (\bar{s}, \bar{f})\bigr)
\bigm| q'(s, f) = q''(\bar{s}, \bar{f})\bigr\}
\]
est l'ensemble des « repères » $(s, a, \bar{f})$, avec trois projections
$\pi_{01}, \pi_{12}, \pi_{02} : E \to R$ ; il se scinde selon l'orientation du
repère en $E^{+}$ et $E^{-}$ --- les triangles positifs et négatifs de la
page 63, dont la page 77 redessine la double pyramide, à côté d'un pavage par
des pentagones où un triangle barycentrique est marqué $s_{0}, s_{1}, s_{2},
d_{01}, d_{12}, d_{20}, e^{+}$.\footnote{Les formules pour $\sigma_{0}, \sigma_{1},
\sigma_{2}$ sur $E$ sont surchargées et leurs indices douteux ; on n'en
restitue que le principe.}

\section*{11. Groupes profinis spéciaux et action de Galois (pages 79 à 90)}

\pagerange{79}{80}
\subsection*{Définition (pages 79 et 80)}

Un \emph{groupe profini spécial} est un système $\Sigma = (T, \Pi, I,
(\Phi_{i})_{i \in I})$ --- $T \simeq \widehat{\mathbb{Z}}$, $\Pi$ profini, $I$
fini, $\Phi_{i} \in \mathrm{Homext}(T, \Pi)$ --- isomorphe à l'un des systèmes
\[
\Sigma_{\nu} : \quad T = \widehat{\mathbb{Z}}, \quad \Pi_{\nu} = \widehat{L}_{\nu},
\quad L_{\nu} = \langle \lambda_{0}, \ldots, \lambda_{\nu} \mid \lambda_{0}\lambda_{1}
\cdots \lambda_{\nu} = 1 \rangle,
\]
\[
I = \{0, \ldots, \nu\}, \qquad \Phi_{i} = [\,1 \mapsto \lambda_{i}\,],
\]
$\nu \geqslant 1$ ; $\Pi_{\nu}$ est profini libre à $\nu$ générateurs. C'est le
modèle profini de la sphère privée de $\nu + 1$ points, avec ses lacets
autour des trous. L'entier $\nu$ vaut $\operatorname{card} I - 1$, et c'est le
rang de $\Pi_{\mathrm{ab}}$. Les $\Sigma$ de $\nu$ donné forment un groupoïde
connexe, équivalent à celui des torseurs sous
\[
\widetilde{A}_{\nu} = \mathrm{Aut}(\Sigma_{\nu}),
\]
qui s'envoie \emph{surjectivement} sur $\mathfrak{S}_{\nu+1} \times
\widehat{\mathbb{Z}}^{*}$ (permutation de $I$, action sur $T$). On note $A_{\nu}$
le noyau vers $\mathfrak{S}_{\nu+1}$ et $A^{\circ}_{\nu}$ le noyau vers le produit.
Les automorphismes intérieurs donnent $\Pi_{\nu} \hookrightarrow A^{\circ}_{\nu}$,
injectif pour $\nu \geqslant 2$, et l'on note $\widetilde{\Gamma}_{\nu}$,
$\Gamma_{\nu}$, $\Gamma^{\circ}_{\nu}$ les quotients de $\widetilde{A}_{\nu}$,
$A_{\nu}$, $A^{\circ}_{\nu}$ par $\Pi_{\nu}$ : ce sont les groupes d'automorphismes
\emph{extérieurs} correspondants. D'où des suites exactes
\[
1 \to \Pi_{\nu} \to A^{\circ}_{\nu} \to \Gamma^{\circ}_{\nu} \to 1, \qquad
1 \to A_{\nu} \to \widetilde{A}_{\nu} \to \mathfrak{S}_{\nu+1} \to 1, \qquad
1 \to A^{\circ}_{\nu} \to A_{\nu} \to \widehat{\mathbb{Z}}^{*} \to 1,
\]
et une filtration de $\widetilde{A}_{\nu}$ de quotients successifs $\Pi_{\nu}$,
$\Gamma^{\circ}_{\nu}$, $\widehat{\mathbb{Z}}^{*}$, $\mathfrak{S}_{\nu+1}$.%
\footnote{Pour $\nu = 1$, $\Pi_{1} = \widehat{\mathbb{Z}}$ est abélien et ses
automorphismes intérieurs sont triviaux : l'injectivité demande $\nu
\geqslant 2$, hypothèse que la page 88 fait. La première suite exacte de la
page a $\Gamma_{\nu}$ sans l'exposant $\circ$ annoncé.}

\emph{Exemple} : $X$ la droite projective sur un corps $k$ algébriquement clos
de caractéristique $0$, $I \subset X(k)$ de cardinal $\nu + 1$, $T = \varprojlim
\mu_{n}(k)$, $\Pi = \pi_{1}(X \smallsetminus I, x)$, et $\Phi_{i}$ les classes
d'inertie.

\pagerange{81}{86}
\subsection*{L'action de Galois (pages 81 à 86)}

Le système $\Sigma(X/I/\xi/k) = (T, I, \pi, (\Phi_{i}))$ de l'exemple ne dépend
que d'un foncteur fibre $\xi$ sur les revêtements étales de $X \smallsetminus
I$, et il en dépend fonctoriellement ; quand $\xi$ varie, il n'est défini qu'à
automorphisme intérieur près, et définit un torseur à droite $P(X/I/k)$ sous
$\widetilde{\Gamma}_{\nu}$, dont le commutant est le groupe des automorphismes
extérieurs de $\Sigma$. Si $X/I/k$ provient d'une situation $X_{0}/I_{0}/k_{0}$
par extension de la base à $k = \bar{k}_{0}$, le groupe $\Gamma =
\mathrm{Gal}(\bar{k}_{0}/k_{0})$ opère par transport de structure, d'où un
homomorphisme
\[
\Gamma \longrightarrow {}^{P(X/I/k)}\widetilde{\Gamma}_{\nu} \simeq
\mathrm{Aut}(\Sigma^{\mathrm{ext}})
\]
vers la forme tordue de $\widetilde{\Gamma}_{\nu}$ par $P$ : c'est la donnée qui
redonne, par image réciproque, l'extension $1 \to \pi_{1}(X \smallsetminus I) \to
\pi_{1}(X_{0} \smallsetminus I_{0}) \to \pi_{1}(k_{0}) \to 1$, et son composé avec
$\widetilde{\Gamma}_{\nu} \to \mathrm{Aut}(T) = \widehat{\mathbb{Z}}^{*}$ est le
caractère cyclotomique.\footnote{Nom nôtre : la « représentation galoisienne
extérieure » du groupe fondamental, qu'étudieront Y.~Ihara, P.~Deligne
(« Le groupe fondamental de la droite projective moins trois points », 1989)
et d'autres. La page dit « homom.\ de Kummer--Tate ».}

\emph{Sur $\mathbb{Q}$, avec $I = \mu_{n}$.} Soient $k_{0} = \mathbb{Q}$, $X_{0} =
\mathbb{P}^{1}_{\mathbb{Q}}$, $n = \nu + 1$ et $I_{0} = \mu_{n}$, dont les points
géométriques sont les $\mu_{n}(\bar{\mathbb{Q}}) \subset \mathbb{C}$. Les lacets
issus de $0$ qui contournent les racines $n$-ièmes de l'unité dans l'ordre
donnent un isomorphisme canonique $\Sigma^{\mathrm{ext}} \simeq
\Sigma_{\nu}^{\mathrm{ext}}$, donc $P$ est trivial et l'on a
\[
i_{\nu} : \Gamma \longrightarrow \widetilde{\Gamma}_{\nu} .
\]
\footnote{Les pages 83 à 86 écrivent $\widetilde{\mathcal{T}}_{\nu}$, un $T$
cursif, pour ce que la page 81 appelle $\widetilde{\Gamma}_{\nu}$ ; on garde la
lettre de la page 81, et l'on réserve $\mathcal{T}_{\nu}$ au groupe discret de
la page 89.} Son composé $\varepsilon_{\nu} i_{\nu} : \Gamma \to
\mathfrak{S}_{n}$ est l'action de $\Gamma$ sur $\mu_{n}$, qui se factorise par
$\chi_{n} : \Gamma \to (\mathbb{Z}/n)^{*} \subset \mathfrak{S}_{n}$.\footnote{La
page 83 écrit $\Gamma \to \mu_{\nu+1} \simeq \mathbb{Z}/\nu{+}1$ ; c'est le
groupe des unités, comme le dit la relation encadrée de la page 84,
$\varepsilon_{\nu}\, i_{\nu} = \chi_{n} \bmod n$.} L'intérêt de ce choix de $I$
est une symétrie : le groupe diédral $(\mathbb{Z}/2) \ltimes \mu_{n}$, rotations
$z \mapsto \zeta z$ et inversion $z \mapsto 1/z$, préserve $(X_{0}, I_{0})$ ;
avec $\Gamma$ il donne
\[
(\Gamma \times \mathbb{Z}/2) \ltimes \mu_{n}(\bar{\mathbb{Q}}) \longrightarrow
\widetilde{\Gamma}_{\nu}.
\]
La rotation induit l'automorphisme $\rho_{\nu}$ de $\widehat{L}_{\nu}$ qui
permute circulairement les générateurs, $\lambda_{i} \mapsto \lambda_{i+1}$ ;
l'inversion, qui fixe $1$ et renverse l'ordre circulaire, induit un
automorphisme $t_{\nu}$ qui envoie $\lambda_{0}$ sur $\lambda_{0}$ et chaque
$\lambda_{i}$ sur un conjugué de $\lambda_{\nu+1-i}$, par exemple
\[
t_{\nu}(\lambda_{i}) = h_{i}\, \lambda_{\nu+1-i}\, h_{i}^{-1}, \qquad
h_{i} = (\lambda_{\nu+2-i} \cdots \lambda_{\nu})^{-1} \quad (1 \leqslant i \leqslant \nu),
\]
qui respecte la relation $\lambda_{0}\lambda_{1} \cdots \lambda_{\nu} = 1$.%
\footnote{La page écrit $t_{\nu}(\lambda_{0}, \ldots, \lambda_{\nu}) =
(\lambda_{0}, g_{\nu}\lambda_{\nu}g_{\nu}^{-1}, \ldots, g_{2}\lambda_{2}g_{2}^{-1},
\lambda_{1})$, $g_{i} = \lambda_{0} \cdots \lambda_{i-1}$, avec en marge
« vérifier ? ». Comme $g_{i}\lambda_{i}g_{i}^{-1} = g_{i+1}g_{i}^{-1}$, le
produit des images vaut $\lambda_{0}\, g_{2}^{-1}\lambda_{1} =
\lambda_{0}\lambda_{1}^{-1}\lambda_{0}^{-1}\lambda_{1}$, qui n'est pas $1$ : la
formule ne définit pas un automorphisme de $L_{\nu}$. La nôtre se vérifie par
le même télescopage. Ni l'une ni l'autre n'est déduite de la géométrie de
$z \mapsto 1/z$, qui déplace le point base $0$ ; seule leur classe dans
$\widetilde{\Gamma}_{\nu}$ importe.} On a alors les relations
\[
\varepsilon_{\nu}\, i_{\nu} = \chi_{n} \bmod n, \qquad
i_{\nu}(\Gamma) \text{ commute à } t_{\nu}, \qquad
i_{\nu}(\gamma)\, \rho_{\nu}(\zeta)\, i_{\nu}(\gamma)^{-1} =
\rho_{\nu}(\gamma(\zeta)) ,
\]
la dernière parce que conjuguer une rotation par Galois change la racine de
l'unité dont elle est la multiplication. Soit $\Gamma^{!}_{\nu} \subset
\widetilde{\Gamma}_{\nu}$ le sous-groupe des éléments qui satisfont ces
relations (pour un caractère $\chi$ convenable) ; on a
\[
1 \to \Gamma^{!\circ}_{\nu} \to \Gamma^{!}_{\nu} \to \widehat{\mathbb{Z}}^{*} \to 1,
\qquad
\Gamma^{!\circ}_{\nu} = \{\gamma \in \Gamma^{\circ}_{\nu} \mid \gamma \text{ commute
aux } \rho_{\nu}(\zeta) \text{ et à } t_{\nu}\},
\]
et $i_{\nu}$ prend ses valeurs dans $\Gamma^{!}_{\nu}$.\footnote{La page note
$\mathcal{T}^{!}_{\nu}$, $\mathcal{T}^{!\,0}_{\nu}$. La surjectivité sur
$\widehat{\mathbb{Z}}^{*}$ vient de celle de $\chi$.}

\emph{Le cas exceptionnel $\nu = 2$, $I = \{0, 1, \infty\}$.} Les trois points sont
alors rationnels et $\mathfrak{S}_{3}$ opère sur $(X_{0}, I_{0})$ en commutant à
$\Gamma$, d'où $\Gamma \times \mathfrak{S}_{3} \to \widetilde{\Gamma}_{2}$. Ici,
exceptionnellement, $\widetilde{\Gamma}_{2} \simeq \mathfrak{S}_{3} \ltimes
\Gamma_{2}$ ; comme $\Gamma$ fixe les trois points, son image dans
$\mathfrak{S}_{3}$ est triviale, et
\[
\boxed{\ \Gamma \longrightarrow \Gamma_{2}^{\mathfrak{S}_{3}}\ }
\]
--- l'action de Galois sur le groupe fondamental profini de la droite privée
de trois points prend ses valeurs dans le centralisateur de $\mathfrak{S}_{3}$.
La page voudrait des isomorphismes canoniques $\Gamma^{!}_{2} \simeq
\Gamma_{2}^{\mathfrak{S}_{3}}$ et $\Gamma^{!}_{2} \simeq \Gamma^{!}_{\nu}$, avec
un point d'interrogation, et s'interrompt.\footnote{Noms nôtres. Que
$\Gamma \to \mathrm{Out}(\widehat{\pi}_{1}(\mathbb{P}^{1} \smallsetminus \{0, 1,
\infty\}))$ soit injectif découle du théorème de G.~V.~Belyi (1979), postérieur
à ces pages, que l'\emph{Esquisse} saluera. Le groupe défini par la
compatibilité à $\mathfrak{S}_{3}$ (relations « à deux et à trois cycles ») et
par une relation venant de $M_{0,5}$ est le groupe de Grothendieck--Teichmüller
$\widehat{GT}$ de V.~Drinfeld (1990), étudié par Y.~Ihara ; le souhait
$\Gamma^{!}_{2} \simeq \Gamma^{!}_{\nu}$ est une forme de l'idée que ce même
groupe agit sur toute la tour des $M_{0,n}$. La page ne contient évidemment
rien de cela ; elle en pose le premier étage.}

\pagerange{88}{90}
\subsection*{Le groupe de Teichmüller discret et son complété (pages 88 à 90)}

La page 88 reprend la définition sous forme d'hypothèse $\mathrm{Hyp}_{\nu}$
($\nu \geqslant 2$) et pose $A_{\nu} = \mathrm{Aut}(\widehat{L}_{\nu}, T_{0},
(\Phi_{0i}))$, les $\Phi_{0i}$ étant fixés un à un : $\chi_{\nu} : A_{\nu} \to
\widehat{\mathbb{Z}}^{*}$ est surjectif, $\widehat{L}_{\nu}$ s'injecte dans
$\mathrm{Ker}\,\chi_{\nu}$ par les automorphismes intérieurs, son image est
distinguée ($u\, \mathrm{int}(g)\, u^{-1} = \mathrm{int}(u(g))$), et l'on veut
élucider $\Gamma_{\nu} = A_{\nu}/\widehat{L}_{\nu}$.\footnote{La page écrit
$\mathrm{Homext}(L, \pi)$ là où l'on attend $T$ ; et « on prouve que c'est
surjectif » sans la preuve.}

Le modèle discret : $B_{\nu} = \mathrm{Aut}(L_{\nu}, \mathbb{Z}, (\Psi_{i}))$, $\Psi_{i}$
la classe de $1 \mapsto \lambda_{i}$, s'envoie sur $\mathbb{Z}^{*} = \{\pm 1\}$, et
\[
\mathcal{T}_{\nu} = B_{\nu}/L_{\nu}, \qquad
1 \to \mathcal{T}^{\circ}_{\nu} \to \mathcal{T}_{\nu} \to \mathbb{Z}/2 \to 1,
\]
est le groupe de Teichmüller discret de la sphère à $\nu + 1$ trous : par
Dehn--Nielsen--Baer, le groupe des classes d'isotopie d'homéomorphismes qui
fixent chaque trou, $\mathcal{T}^{\circ}_{\nu}$ correspondant à ceux qui
respectent l'orientation. On a $\mathcal{T}^{\circ}_{2} = 1$, et pour $\nu
\geqslant 3$, $\mathcal{T}^{\circ}_{\nu}$ est non trivial : c'est le groupe des
tresses sphériques colorées à $\nu + 1$ brins modulo son centre d'ordre $2$,
et le groupe fondamental de l'espace des configurations de $\nu + 1$ points
distincts de la sphère modulo $PGL_{2}$.\footnote{Noms nôtres : groupe modulaire
pur $\mathrm{PMod}(S_{0,\nu+1}) \simeq \pi_{1}(M_{0,\nu+1})$. La page dit « groupe
des tresses sphériques colorées » sans le quotient par le centre, et demande
« $[\simeq \pi_{1}(\widetilde{D}_{\nu+1})\ ?]$ ».} La complétion donne un diagramme
\[
\widehat{L}_{\nu} \hookrightarrow \widehat{B}_{\nu} \to A_{\nu}, \qquad
\widehat{\mathcal{T}}_{\nu} = \widehat{B}_{\nu}/\widehat{L}_{\nu} \longrightarrow
\Gamma_{\nu} = A_{\nu}/\widehat{L}_{\nu}, \qquad
\widehat{\mathcal{T}}^{\circ}_{\nu} \to \Gamma^{\circ}_{\nu},
\]
compatible aux flèches vers $\mathbb{Z}/2 \subset \widehat{\mathbb{Z}}^{*}$, et
trois questions : $\widehat{\mathcal{T}}_{\nu}$ est-il distingué dans
$\Gamma_{\nu}$ ? $\widehat{\mathcal{T}}^{\circ}_{\nu}$ dans $\Gamma^{\circ}_{\nu}$ ?
Quand a-t-on $\widehat{\mathcal{T}}^{\circ}_{\nu} = \Gamma^{\circ}_{\nu}$ ? ---
jamais pour $\nu = 2$, où le premier est trivial et le second contient
l'image de Galois.\footnote{La page dessine $\widehat{\mathcal{T}}_{\nu} \subset
\Gamma_{\nu}$ comme une inclusion. L'injectivité de $\widehat{B}_{\nu} \to A_{\nu}$
est la « propriété des sous-groupes de congruence » en genre $0$, démontrée par
M.~Asada (2001). La dernière entrée de la ligne du bas du premier diagramme
est écrite $\mathbb{Z}^{*}$ sur la page, pour $\widehat{\mathbb{Z}}^{*}$.}

\emph{Variante} (page 90) : pour $\varepsilon \in [1, \infty]^{\nu+1}$, le groupe
$L_{\nu}^{\varepsilon} = \langle \lambda_{0}, \ldots, \lambda_{\nu} \mid \lambda_{0}
\cdots \lambda_{\nu} = \lambda_{0}^{\varepsilon_{0}} = \cdots =
\lambda_{\nu}^{\varepsilon_{\nu}} = 1\rangle$ ; la page s'arrête sur sa
définition.\footnote{Nom nôtre : groupe fondamental orbifold de la sphère à
points coniques, groupe triangulaire pour $\nu = 2$.}

\pagerange{92}{95}
\section*{12. Le torseur des clôtures algébriques (pages 92 à 95)}

Ces pages, d'une autre encre, rendent la construction indépendante du choix
de $\bar{\mathbb{Q}}$. Le cadre abstrait : un groupe $\Gamma$ et un
$\Gamma$-torseur à droite $\mathcal{C}$ ; un $\Gamma$-groupe $\pi$ ; pour $i \in
I$, des $(\pi, \Gamma)$-torseurs $T_{i}$, leurs commutants $\pi_{i}$, et des
sous-$\Gamma$-groupes $I_{i} \subset \pi_{i}$. Tout se tord par $\mathcal{C}$ :
\[
\widetilde{\pi} = \mathcal{C} \wedge^{\Gamma} \pi, \qquad
(\widetilde{\pi}_{i}, \widetilde{T}_{i}, \widetilde{\pi}) = \mathcal{C}
\wedge^{\Gamma} (\pi_{i}, T_{i}, \pi), \qquad
\widetilde{I}_{i} = \mathcal{C} \wedge^{\Gamma} I_{i},
\]
et un élément $\varphi \in \mathcal{C}$ donne $\beta(\varphi) : \pi
\xrightarrow{\sim} \widetilde{\pi}$, avec $\beta(\varphi\gamma) = \beta(\varphi)
\circ \gamma_{\pi}$. L'interprétation : $\Gamma = \mathrm{Aut}(\bar{\mathbb{Q}})$,
$\mathcal{C}$ l'ensemble des isomorphismes de $\bar{\mathbb{Q}}$ avec une
clôture algébrique quelconque, $\pi = \pi_{1}(\mathbb{P}^{1}_{\bar{\mathbb{Q}}}
\smallsetminus I)$ pour $I \subset \mathbb{P}^{1}(\mathbb{Q})$ fini, $T_{i}$ le
torseur des chemins du point base vers un point base au voisinage de $i$,
$\pi_{i}$ son groupe, et $I_{i} \simeq \mathrm{T}$ le sous-groupe d'inertie,
procyclique. Le calcul des pages 94--95 vérifie que les
$\alpha'_{i}(\varphi) = \beta(\varphi) \alpha_{i}(\varphi) : I_{i} \to
\widetilde{\pi}$ satisfont $\alpha'_{i}(\varphi\gamma) = \alpha'_{i}(\varphi)\,
\gamma_{\pi_{i}}$, donc que les sous-groupes d'inertie de $\widetilde{\pi}$ ne
dépendent pas de $\varphi$, et que l'action sur $I_{i}$ se fait par un
caractère $\chi : \Gamma \to \widehat{\mathbb{Z}}^{*}$ indépendant de $i$, dont le
quotient effectif est $\mathrm{Aut}(I_{i}) \simeq
\widehat{\mathbb{Z}}^{*}$.\footnote{Les pages 94--95 écrivent $\pi$ là où l'on
attend $\Gamma$ dans « $\pi \xrightarrow{\chi} \widehat{\mathbb{Z}}^{*}$ »,
« $\mathcal{C} \wedge_{\pi} I_{i}$ » et « $\pi' = \pi/U_{i}$ » ; on lit $\Gamma$,
comme à la page 92. Le dernier terme, « $\mathcal{C} \simeq \varprojlim
\mu_{n}(\bar{\mathbb{Q}})^{*}$ », se lit comme le torseur des générateurs de
$\mathrm{T}$ ; l'ensemble de la page est elliptique.}

\section*{13. Précartes et cartes co-épinglées (pages 96 à 111)}

\pagerange{97}{98}
\subsection*{I) Précartes, et la catégorie des cartes co-épinglées
(pages 97 et 98)}

Une \emph{précarte} est un système $C_{0} = (S, A, F ; p, q, r)$ d'ensembles et
d'applications $p : S \to \mathbb{N}^{*}$, $q : A \to \{1, 2\}$, $r : F \to
\mathbb{N}^{*}$ : la liste des sommets, des arêtes et des faces avec leurs
degrés, une arête ayant un ou deux brins. Les isomorphismes de précartes sont
les triples de bijections qui respectent les degrés.\footnote{La page définit
des morphismes $\varphi$ plus généraux, avec $p(s') = \nu(\varphi, s)\, p(s)$,
$\nu(\varphi, s) = \operatorname{Card} \varphi_{s}^{-1}(s)$, après des conditions
de divisibilité biffées ; seuls les isomorphismes ($\nu = 1$) servent dans la
suite, et la ligne qui l'annonce (« pour simplement iso ») est d'une lecture
incertaine.} Une carte combinatoire finie connexe $C$ a sa précarte
$\mathrm{Fac}(C)$, et le foncteur $\mathrm{Fac}$ est conservatif. Une \emph{carte
co-épinglée} par $C_{0}$ est une carte $C$ munie d'un isomorphisme $u : C_{0}
\xrightarrow{\sim} \mathrm{Fac}(C)$ ; les morphismes sont les isomorphismes de
cartes compatibles aux $u$.

Si $C_{0}$ est fini, cette catégorie est finie, et vide sauf si $\sum_{s} p(s) =
\sum_{a} q(a) = \sum_{f} r(f) = \delta$, le nombre de brins (la
« complexité »). Elle est \emph{rigide} pour les cartes de genre $0$, mais pas
en général.\footnote{La page affirme la rigidité sans restriction (p.~98, et de
nouveau p.~99, avec en marge « vérifier »). Elle est fausse en genre $1$ : la
carte du tore à un sommet, deux arêtes et une face (le carré dont on identifie
les côtés opposés) a pour automorphisme la symétrie $-1$, qui fixe le sommet,
chaque arête et la face, donc chaque co-épinglage. En genre $0$ elle est vraie :
un automorphisme d'ordre fini qui fixe chaque cellule fixe au moins trois
points de la sphère ($V + F = E + 2 \geqslant 3$), et une homographie qui fixe trois
points est l'identité.} Lorsqu'elle l'est, le nombre de classes d'isomorphie
de cartes co-épinglées de type $C_{0}$ est
\[
\sum_{C} \frac{|\mathrm{Aut}(C_{0})|}{|\mathrm{Aut}(C)|}, \qquad
|\mathrm{Aut}(C_{0})| = \prod_{n \in \mathbb{N}^{*}} a_{n}!\, b_{n}!\, c_{n}!,
\]
la somme portant sur les classes de cartes $C$ telles que $\mathrm{Fac}(C)
\simeq C_{0}$, et $a_{n}, b_{n}, c_{n}$ comptant les sommets, arêtes et faces de
degré $n$.\footnote{La page dit : le nombre des classes de cartes de type
$C_{0}$, « multiplié par » $|\mathrm{Aut}(C_{0})|$. C'est juste si chaque carte
est comptée avec le poids $1/|\mathrm{Aut}(C)|$ (cardinal du groupoïde), et
faux sinon : pour l'étoile à trois branches de la page 103, $|\mathrm{Aut}(C_{0})|
= 6$, $\mathrm{Aut}(C) = \mathbb{Z}/3$, et il y a $6/3 = 2$ cartes co-épinglées,
comme le dit la page 103 elle-même.} Le groupe $\gamma_{C_{0}} =
\mathrm{Aut}(C_{0})$ opère sur l'ensemble $\mathcal{C}^{\alpha}_{C_{0}}$ des
classes de cartes co-épinglées de type $(C_{0}, \alpha)$, et le quotient est
l'ensemble des classes de cartes de ce type.

\pagerange{99}{102}
\subsection*{II--III) Le schéma de modules, et l'action de Galois (pages 99 à 102)}

Soit $C_{0}$ fini et $T$ un schéma de base. Un \emph{revêtement co-épinglé} de
$\mathbb{P}^{1}_{T}$ est un revêtement fini $f : X \to \mathbb{P}^{1}_{T}$, étale
au-dessus de $\mathbb{P}^{1}_{T} \smallsetminus \{0, 1, \infty\}$, modérément
ramifié fibre à fibre, muni d'isomorphismes
\[
f^{-1}(0_{T})_{\mathrm{réd}} \simeq S_{T}, \qquad
f^{-1}(1_{T})_{\mathrm{réd}} \simeq A_{T}, \qquad
f^{-1}(\infty_{T})_{\mathrm{réd}} \simeq F_{T}
\]
compatibles aux degrés de ramification. Soit $F_{C_{0}}(T)$ l'ensemble de leurs
classes. Là où la catégorie est rigide, la page affirme que le foncteur
$F_{C_{0}}$ est représentable par un schéma $M_{C_{0}}$ quasi-fini, étale et
séparé sur $\mathbb{Z}$, fini au-dessus de $\mathbb{Z}[1/\delta!]$. En particulier
$(M_{C_{0}})_{\mathbb{Q}}$ est fini étale sur $\mathbb{Q}$, non ramifié en les $p >
\delta$.\footnote{Nom nôtre : c'est un schéma de Hurwitz au sens de W.~Fulton
(« Hurwitz schemes and irreducibility of moduli of algebraic curves »,
\emph{Ann.\ of Math.}, 1969), avec des points de ramification fixés ; la bonne
réduction en $p > \delta$ vient de ce que le groupe de monodromie, contenu dans
$\mathfrak{S}_{\delta}$, est alors d'ordre premier à $p$. La page n'en donne pas
de démonstration. Hors du cas rigide, le foncteur a des automorphismes et
n'est pas représentable ; voir la note sur la page 98.} Le groupe
$\gamma_{C_{0}}$ opère sur $M_{C_{0}}$ ; le stabilisateur d'un point correspond
aux automorphismes du revêtement $f$, indépendamment de l'épinglage, et c'est
le groupe de Galois de $f$ quand $f$ est galoisien.

Pour une classe $\alpha$ de groupes finis, $M^{\alpha}_{C_{0}} \subset
M_{C_{0}}$ est le sous-schéma ouvert et fermé, stable par $\gamma_{C_{0}}$, des
$f$ dont la clôture galoisienne a un groupe de type $\alpha$. Il est vide sauf
si l'ordre $N$ de ce groupe est un multiple commun des degrés et de $\delta$, et
il est non vide si et seulement s'il existe un groupe $G$ de type $\alpha$
engendré par $\rho_{s}, \rho_{f}, \sigma$ avec
\[
\rho_{s}\, \rho_{f}\, \sigma = 1, \qquad \sigma^{2} = 1,
\]
et un sous-groupe $H \subset G$ dont l'intersection des conjugués est triviale,
tels que $C_{0}$ soit la précarte de la carte orientée $G/H$. La page ajoute
que, lorsque $N = \delta$ ($H = 1$), le type de carte orientée n'est pas
déterminé par le seul revêtement.\footnote{Nous rapportons cette remarque sans
pouvoir la reconstituer : pour $H = 1$, le revêtement galoisien détermine son
groupe d'automorphismes et la monodromie à conjugaison simultanée près, donc,
semble-t-il, la carte orientée $G$. La phrase de la page commence par des mots
incertains (« dans un cas qui arrivât »), et c'est peut-être autre chose qui
est visé. La condition de la page sur $N$ et $\delta$ est lue telle quelle.}

\emph{L'action de Galois.} Par le théorème d'existence de Riemann,
\[
M^{\alpha}_{C_{0}}(\bar{\mathbb{Q}}) \simeq M^{\alpha}_{C_{0}}(\mathbb{C}) \simeq
\mathrm{Hom}(R^{\alpha}_{C_{0}}, \mathbb{C}) \simeq \mathcal{C}^{\alpha}_{C_{0}},
\]
où $R^{\alpha}_{C_{0}}$ est l'algèbre finie étale sur $\mathbb{Q}$ qui représente
$(M^{\alpha}_{C_{0}})_{\mathbb{Q}}$ : l'ensemble \emph{combinatoire} des cartes
co-épinglées est l'ensemble des points algébriques d'un $\mathbb{Q}$-schéma fini.
Le groupe $\Gamma = \mathrm{Aut}(\bar{\mathbb{Q}})$ y opère donc, en commutant à
$\gamma_{C_{0}}$. La conjugaison complexe change une carte en son image miroir,
de sorte que l'action de $\Gamma$ sur $\mathcal{C}^{\alpha}_{C_{0}}/\gamma_{C_{0}}$
n'est pas triviale dès qu'une carte du type est chirale ; on va déterminer
l'action de $\gamma_{C_{0}} \times \Gamma$ dans des cas élémentaires.%
\footnote{La page dit que l'action « ne peut pas être triviale » ; elle l'est
pour les types qui n'ont qu'une carte, comme les premiers exemples de la page
103. Le théorème d'existence de Riemann, sous la forme algébrique utilisée ici, est
SGA~1, exposé XII.}

\pagerange{103}{104}
\subsection*{Premiers exemples (pages 103 et 104)}

On note un type par les degrés des sommets, des arêtes et des faces ; les
fonctions $f$ sont normalisées par $f^{-1}(0) =$ sommets, $f^{-1}(1) =$ milieux
d'arêtes, $f^{-1}(\infty) =$ centres des faces.

\begin{itemize}
\item Le segment à une extrémité : $f(z) = z$, $G = 1$.
\item Le segment de $-1$ à $1$ par le sommet $0$ : $f(z) = z^{2}$, $G \simeq
\mathbb{Z}/2$, une seule carte co-épinglée.
\item Le segment à deux sommets $0$ et $1$ : $f(z) = -4z(z-1)$, le milieu
d'arête en $\tfrac12$, $G \simeq \mathbb{Z}/2$, une seule carte.
\item La boucle issue de $0$ : $f(z) = \tfrac14\, z^{2}/(z-1)$, le milieu d'arête
en $z = 2$ (où $f' = 0$ et $f = 1$), une seule carte.
\item L'étoile à trois branches : $f(z) = z^{3}$, $G = \mathbb{Z}/3 \subset
\gamma_{C_{0}} \simeq \mathfrak{S}_{3}$ ; deux cartes co-épinglées, permutées
transitivement par $\gamma_{C_{0}}$, et $\Gamma$ y opère par le caractère
cyclotomique $\chi_{3}$.
\item Le segment prolongé d'une boucle, de type $(3 ; 2, 1 ; 2, 1)$ :
$f(z) = \tfrac{4}{27}\, z^{3}/(z-1)$, avec $f(z) - 1 = (2z-3)^{2}(z+3)/27(z-1)$,
milieux d'arêtes en $\tfrac32$ et $-3$ ; $G \simeq \mathfrak{S}_{3}$,
$\gamma_{C_{0}} = 1$, une seule carte. La page 104 décrit la figure barycentrique
complète comme réunion de l'axe réel et de la courbe $2x(x^{2} + y^{2}) +
(-3x^{2} + y^{2}) = 0$.
\item Son dual : $f(z) = \tfrac{27}{4}\, z(z-1)^{2}$, milieux d'arêtes en
$\tfrac13$ et $\tfrac43$ ; $G = \mathfrak{S}_{3}$, $\gamma_{C_{0}} = 1$, une seule
carte.
\end{itemize}

\pagerange{105}{111}
\subsection*{Un arbre de degré 6 (pages 105 à 111)}

Les pages 105--111, écrites au dos de documents dactylographiés, traitent un
cas où l'action de Galois devrait se voir. Deux arbres plans (I) et (II), à
trois sommets $a, b, c$ de degrés $3, 2, 1$, deux arêtes complètes $u, v$ et deux
arêtes pendantes $x, y$, sont codés par les permutations $\rho_{s}$ (rotation
autour des sommets), $\sigma$ (retournement des arêtes) et $\rho_{f} =
\rho_{s}^{-1}\sigma$ de leurs six brins, que la page écrit en entier. Le
polynôme correspondant est
\[
f(z) = \lambda\, z^{3}(z-1)^{2}(z - c),
\]
un polynôme de degré $6$ (un seul pôle, d'ordre $6$, en $\infty$) avec $a = 0$,
$b = 1$. On calcule
\[
f'(z) = \lambda z^{2}(z-1)\bigl[6z^{2} - (5c+4)z + 3c\bigr] ;
\]
les milieux d'arêtes complètes $u, v$ sont les racines du trinôme, et il faut
$f(u) = f(v) = 1$. Réduisant $u^{3}(u-1)^{2}(u-c)$ modulo $u^{2} = \alpha u +
\beta$, $\alpha = \tfrac{5c}{6} + \tfrac23$, $\beta = -\tfrac{c}{2}$, on obtient
$f(u)/\lambda = P(c)\, u + Q(c)$, d'où, puisque $u \neq v$,
\[
P(c) = 0, \qquad \lambda = \frac{1}{Q(c)}, \qquad c \neq 0, 1, \qquad
\delta(c) = (5c+4)^{2} - 72c = 25c^{2} - 32c + 16 \neq 0,
\]
la dernière condition exprimant $u \neq v$. La page établit des pièces de
$P$ et $Q$ --- par exemple $(\alpha - 2)(\alpha - c) + \beta + 1 =
-\tfrac{1}{36}(5c^{2} - 10c - 4)$ et $\alpha(\alpha^{2} + 2\beta) =
\tfrac{1}{6^{3}}(5c+4)(25c^{2} + 4c + 16)$ ---, sans écrire $P$ en entier. Elle
montre que les zéros de $P$ ne sont ni $0$, ni $1$, ni des zéros de $Q$ :
$P(c) = Q(c) = 0$ forcerait $f(u) = 0$, donc $u \in \{0, 1, c\}$, et chacun de
ces cas donne $c \in \{0, 1\}$. Elle demande enfin si ces zéros pourraient
annuler $\delta$, dont les racines sont $\tfrac{16}{25} \pm \tfrac{12}{25}\, i$,
et la question reste ouverte.\footnote{La page écrit $c = \tfrac{16}{25} \pm
\tfrac{24}{25}\, i$ ; le discriminant de $25c^{2} - 32c + 16$ est $-576 =
-24^{2}$, et les racines sont $(32 \pm 24i)/50$. Nom nôtre pour ces
polynômes : « polynômes de Chebyshev généralisés » ou de Shabat (G.~Shabat
et V.~Voevodsky, 1990). Les deux arbres (I) et (II) ont même type, et chacun
donne une valeur de $c$ ; que le calcul vise à savoir s'ils sont conjugués sous
Galois est une lecture nôtre. La page 109 note aussi la relation $2(\sum p + \sum q + \sum r) = 6N$ entre
degrés et ordre du groupe, et la page 111 commence une approche par
$f'/f$.}

\pagerange{113}{119}
\section*{14. Cartes sphériques en bas degrés (pages 112 à 119)}

Sous ce titre, les pages 113 à 119 recensent à la main les cartes connexes de
la sphère par nombre de brins (le « degré » $n = \sum_{s} q_{s} = \sum_{f}
p_{f}$), de $1$ à $6$, classées par la liste des degrés des sommets. Chaque
carte est dessinée, avec son groupe d'automorphismes --- symétries renversant
l'orientation comprises, comme le montrent les valeurs $\mathbb{Z}/2$ pour un
segment et $D_{n}$ pour l'étoile à $n$ branches --- et son type (degrés des
sommets, des arêtes, des faces). Les cartes duales, qui échangent sommets et
faces, sont réunies par une accolade.\footnote{Lecture nôtre de l'accolade : les
types qu'elle relie sont, dans tous les cas lisibles, obtenus l'un de l'autre
en échangeant la première et la dernière entrée.}

Les cas impossibles sont relevés : quatre sommets de degrés $1, 1, 1, 2$ en
degré $5$, et cinq sommets de degrés $2, 1, 1, 1, 1$ en degré $6$, car une carte
connexe à $V$ sommets a au moins $V - 1$ arêtes complètes, soit $2V - 2$ brins.
Le mot « rigide » marque certaines cartes, la première
fois en degré $4$. En degré $6$, cinq types notés $(\alpha)$ à
$(\eta)$, d'une autre encre, renvoient à des cartes de la page 118. Le recensement
est celui de la page ; il n'a pas été vérifié ici carte par carte, et la page
116 elle-même relève un type écrit deux fois.\footnote{Plusieurs chiffres des
types sont en surcharge ou sous une tache (p.~116, p.~119), et plusieurs
dessins biffés ; voir la transcription.}

\section*{15. « Opérations » sphériques sur les cartes (pages 120 à 131)}

\pagerange{121}{127}
\subsection*{Ce qu'est une opération (pages 121 à 127)}

Le titre de la couverture (page 120) annonce une « axiomatique provisoire ».
Soit $\mathcal{F}$ un préfaisceau sur la catégorie des cartes finies orientées
(morphismes : les revêtements ramifiés), qui transforme sommes en produits et
tel que $\mathcal{F}(X/H) = \mathcal{F}(X)^{H}$ pour un groupe fini $H$ opérant
avec ramification isolée. $\mathrm{Aut}(C)$ opère sur $\mathcal{F}(X)$, et l'on
pose
\[
\mathcal{F}^{\natural}(C) = \mathcal{F}(X)/\mathrm{Aut}^{\circ}(C),
\]
quotient par les automorphismes isotopes à l'identité ; c'est encore un
préfaisceau, et il s'étend aux cartes munies d'un groupe fini par
$\mathcal{F}^{\natural}(C, G) = \mathcal{F}(C)^{G}/\mathrm{Aut}(C,
G)^{\circ}$.\footnote{L'exposant de $\mathcal{F}$ est un petit signe crochu,
rendu $\natural$ dans la transcription ; « préfaisceau » est lu avec un
doute.} Le programme : trouver toutes les solutions $\mathcal{F}^{\natural}$,
foncteurs des cartes finies à isomorphisme près vers les ensembles finis. Un
argument formel donne
\[
\Gamma(\mathcal{F}^{\natural}) \simeq \varinjlim_{H_{\alpha}}
\mathcal{F}^{\natural}(C_{\alpha}, G_{\alpha}),
\]
où $H_{\alpha}$ parcourt les sous-groupes distingués libres d'indice fini du
groupe cartographique orienté $\mathcal{C}_{2}^{+}$, $C_{\alpha}$ est la carte
régulière associée et $G_{\alpha} = \mathcal{C}_{2}^{+}/H_{\alpha}$.\footnote{La
page dit « sous-groupes libres d'indice fini » ; pour que
$\mathcal{C}_{2}^{+}/H_{\alpha}$ soit un groupe, il les faut distingués. Le
groupe $\mathcal{C}_{2}^{+} = \langle \rho_{s}, \rho_{f}, \sigma \mid \sigma^{2} =
1, \rho_{s}\rho_{f}\sigma = 1\rangle \simeq \mathbb{Z} * \mathbb{Z}/2$ est le
« groupe cartographique orienté » de l'\emph{Esquisse} ; le nom est déjà sur
la page.}

\emph{Exemple.} Pour $\mathcal{F} = \mathfrak{P}$ (l'ensemble des parties),
\[
\Gamma(\mathcal{F}^{\natural}) \simeq \mathfrak{P}(S)/\mathrm{Aut}^{\circ}(S),
\]
où $S$ est la sphère de référence munie du segment $[0, 1]$ et du point
$\infty$ : une « opération » est une partie $A$ de la sphère de référence, à
isotopie près, et elle agit sur une carte $C = (X, K, S)$ en prenant
$A_{X} = f^{-1}(A)$ pour le revêtement $f : X \to S$ qui définit
$C$.\footnote{En théorie non orientée, on trouverait le quotient par le groupe
des automorphismes qui commutent à la conjugaison complexe $\underline{\sigma}$,
symétrie par rapport à l'équateur qui passe par $0, 1, \infty$ ; exposant
souligné lu avec un doute.} Les $A_{X}$ sont des sous-1-complexes si et
seulement si $A$ en est un ; $A_{X} \cup K$ en est un si et seulement si $A \cup
[0, 1]$ est un 1-complexe ; et les $A_{X}$ définissent de nouvelles cartes si et
seulement si a) $A$ est connexe et b) $A$ contient au plus un des points $0, 1,
\infty$.\footnote{Le passage entre a) et b) est un palimpseste (p.~124), et un
ajout encadré, dont la place est incertaine, précise a) : « plus des points $0,
1, \infty$ qui sont $\notin A$, i.e.\ définit une carte sur $S$ ».} On choisit
alors sommets, centres des faces et milieux d'arêtes de $A$ selon trois règles
(page 124). À l'exclusion d'un cas exceptionnel, $A$ avec ces choix est
$\varphi^{-1}([0, 1])$ pour une application $\varphi : S \to S$ ramifiée au-dessus
des seuls points $0, 1, \infty$, avec ramification $2$ au-dessus de $1$, et
$A_{X} = (\varphi f)^{-1}([0, 1])$ : composer avec $\varphi$ ne fait pas sortir
des revêtements ramifiés au-dessus de $0, 1, \infty$, puisque $\varphi^{-1}\{0, 1,
\infty\} \supset \{0, 1, \infty\}$. « Que se passe-t-il dans le cas
exceptionnel ? » --- la page 127 s'arrête sur cette question.\footnote{Les
pages se suivent dans l'ordre 122, 124 et 123, 127 : la page 124 continue la
dernière phrase de la page 122, la page 127 la dernière de la page 123, dont
le début est ailleurs. La page 126 est un feuillet de calculs : les
homographies $z \mapsto z/(cz + 1 - c)$ qui fixent $0$ et $1$ sont involutives
pour $c = 0$ ou $2$, et $g(z) = (2f(z) - 1)^{2}$ replie $[0, 1]$ sur lui-même
($g = 0 \iff f = \tfrac12$, $g = 1 \iff f \in \{0, 1\}$).}

\pagerange{125}{125}
\subsection*{Un théorème de rigidification (page 125)}

\textbf{Théorème.} \emph{Soient $C = (X, K, S)$ une carte connexe et $G$ un groupe
opérant fidèlement sur $X$ par automorphismes de $C$. Conditions
équivalentes :} a) \emph{$G \to \mathrm{Aut}_{\mathrm{isot}}(C)$ est injectif : un
$g$ isotope à l'identité est l'identité ;} b) \emph{les stabilisateurs sont
finis ;} c) \emph{il existe une rigidification de $X$ invariante par $G$ ;} d)
\emph{$(C, G)$ est la réalisation topologique d'une action de $G$ sur une carte
combinatoire. De plus l'espace des rigidifications invariantes par $G$ est
contractile.}\footnote{La numérotation des conditions a été remaniée sur la
page, et la condition b) est un ajout de lecture incertaine ; on suit la
dernière main, sous réserve. La page dit « a) $\Rightarrow$ b) $\Rightarrow$ c)
évident » ; l'ordre des conditions ayant changé, on ne sait pas avec certitude
à quelles conditions ces lettres renvoient.}

La page démontre b) $\Rightarrow$ a) en le ramenant à : \emph{un automorphisme
d'ordre fini d'une carte, isotope à l'identité, est l'identité} ; et cela à un
disque fermé : un homéomorphisme d'ordre fini du disque qui est l'identité sur
le bord est l'identité. La page s'arrête au milieu de cet argument.%
\footnote{Ce dernier énoncé est classique (un homéomorphisme d'ordre fini du
disque est conjugué à une rotation : B.~de Kerékjártó, S.~Eilenberg) ; la
référence est nôtre. On notera que « isotope à l'identité » est plus fort que
« fixant chaque cellule » : la symétrie $-1$ du tore de la note sur la page 98
fixe chaque cellule sans être isotope à l'identité.}

\pagerange{128}{130}
\subsection*{Cartes holomorphes (pages 128 à 130)}

Soit $X$ une surface de Riemann et $C = (K, S)$ une carte sur $X$. La carte
est \emph{holomorphe} s'il existe $\varphi : X \to \mathbb{P}^{1}_{\mathbb{C}}$
holomorphe, étale au-dessus de $\mathbb{P}^{1} \smallsetminus \{0, 1, \infty\}$,
d'indice de ramification $2$ en tous les points au-dessus de $1$, avec $(K, S) =
\varphi^{-1}([0, 1], 0)$.\footnote{Nom nôtre : une fonction de Belyi « propre »
(\emph{clean}), dont le dessin est l'image réciproque du segment.} Une telle
$\varphi$ est-elle unique ? Deux solutions diffèrent par un automorphisme
holomorphe de la carte isotope à l'identité, et la question est de savoir s'il
en existe d'autres que l'identité. Il en existe pour la sphère partagée par un
équateur portant un ou deux sommets :

\begin{itemize}
\item[I)] un sommet et une arête (le sommet en $\infty$, l'arête la droite
réelle) : le groupe $z \mapsto az + b$, $a > 0$, $b \in \mathbb{R}$, de dimension
$2$ ;
\item[II)] deux sommets et deux arêtes (les sommets $0$ et $\infty$) : le groupe
$z \mapsto az$, $a > 0$, de dimension $1$ ;
\item[III)] la carte duale de I), une arête et deux sommets.
\end{itemize}

\footnote{La page donne pour I) « $\mathrm{Aff}(1, \mathbb{R})$ », second argument
surchargé ; pour II) « $\mathbb{C}^{*}$ de dim.\ $1$ » --- une homothétie
complexe ne préserve la droite réelle et chacune des deux demi-droites que si
son rapport est réel positif ; pour III) « $\mathrm{Aff}(1, \mathbb{C})$ », que
nous ne savons pas concilier avec la description : une homographie qui
respecte un segment et ses deux extrémités est réelle, et ces homographies
forment un groupe de dimension $1$. Le groupe de III) reste douteux.}

Si l'on exige que chaque face ait au moins trois sommets, gagne-t-on
l'unicité ? Grothendieck le croit, « grâce à la théorie de la représentation
conforme », au moins pour les cartes finies. Il affirme ensuite que les seuls
autres cas à écarter sont les cartes de l'« orange » d'ordre $n$ (deux pôles,
$n$ méridiens), et que pour $n \geqslant 2$ celles-ci n'ont pas d'automorphisme
holomorphe non trivial. Ce dernier point est faux : $z \mapsto az$, $a > 0$,
préserve chaque méridien et chaque fuseau de l'orange de pôles $0$ et $\infty$,
quel que soit $n$ --- le cas II) est l'orange d'ordre $2$. La conclusion de la
page subsiste néanmoins : \emph{si l'on se donne de plus $\varphi^{-1}(\infty)$,
les centres des faces, les cas I), II), III) et les oranges sont éliminés} ; si
l'on se donne seulement $\varphi^{-1}(1)$, les milieux d'arêtes, II), III) et les
oranges le sont, mais non I). Il reste à « tirer au clair entièrement
l'argument de représentation conforme, et voir s'il s'applique aux cartes
infinies ».

\pagerange{131}{131}
\subsection*{Opérations holomorphes (page 131)}

Une \emph{opération} orientée $\Omega$ est définie par une carte $(K', S')$ sur
la sphère de référence $S$, à isotopie relative à $[0, 1] \cup \{\infty\}$ près.
La page compare des conditions : a) pour toute carte holomorphe $C$, $\Omega(C)$
a un représentant compatible à la structure complexe ; b) les réalisations
holomorphes canoniques de $C$ et de $\Omega(C)$ sont des surfaces isomorphes ;
c) pour toute carte $C$ il existe une structure complexe compatible à la fois
à $C$ et à $\Omega(C)$ ; a'), b'), c') les mêmes avec groupe d'automorphismes ;
et d) il existe $\varphi : S \to S$ holomorphe « spéciale » --- étale au-dessus
de $S \smallsetminus \{0, 1, \infty\}$, de ramification au plus $2$ au-dessus
de $1$, envoyant $\{0, \infty\}$ dans $\{0, \infty\}$ et fixant $1$ --- telle
que $(K', S')$ soit isotope à $(\varphi^{-1}([0, 1]), \varphi^{-1}(0))$ relativement
à $[0, 1] \cup \{\infty\}$. La page s'arrête sur cette condition, sans avoir
dit lesquelles sont équivalentes ; d) est la forme que prend l'exemple de la
page 123, et c'est elle qui rend une opération algébrique : $(\varphi
f)^{-1}([0, 1])$ est encore le dessin d'une fonction ramifiée au-dessus des
seuls points $0, 1, \infty$.\footnote{La dernière phrase est de nous : c'est ce
que les pages 123 et 127 préparent, mais la page 131 ne le dit pas. Le « $\{0\}$ » de
la condition d) est lu $\varphi^{-1}(0)$ ; « bout arrondi » et plusieurs mots de
la page sont incertains.}

\end{document}
